\documentclass[11pt]{article}
\usepackage[margin=1in]{geometry}
\usepackage[T1]{fontenc}
\usepackage[utf8]{inputenc}
\usepackage[french]{babel}
\usepackage{amsmath,amssymb,amsthm}
\usepackage[colorlinks=true,linkcolor=blue,urlcolor=blue]{hyperref}
\newtheorem{problem}{Problème}
\newenvironment{refs}{\par\small\noindent\emph{Références :}\begin{itemize}\setlength\itemsep{0pt}}{\end{itemize}\normalsize}
\title{Hors de la Caverne \\ \Large Statistique}
\author{}
\date{}
\begin{document}
\maketitle
\noindent\emph{Des problèmes, pas de réponses.} Chaque problème ci-dessous est posé
sans solution. Les références indiquent où trouver les connaissances nécessaires.
\medskip


\section*{Collège}

\begin{problem}[{Le salaire moyen chez Duval --- Collège}]
\textbf{STA-41.} \textbf{Problème.} L'entreprise Duval emploie dix personnes : huit ouvriers payés 1 500 € par
mois, un chef d'atelier payé 2 400 € et le patron, qui se verse 21 600 €.
\begin{enumerate}
\item[(a)] Calculez la moyenne et la médiane de ces dix salaires. Justifiez chaque calcul, et en
particulier la façon dont vous déterminez la médiane d'une série de dix valeurs.
\item[(b)] L'entreprise affiche sur son site le salaire moyen de ses employés. Expliquez pourquoi cette
information, bien qu'exacte, donne une image fausse de l'entreprise. Dites quel indicateur vous
publieriez à sa place et défendez votre choix.
\item[(c)] Le patron double son propre salaire, les neuf autres restant inchangés. Dites ce que
deviennent la moyenne, la médiane et l'étendue, en justifiant les trois réponses. Énoncez ensuite,
en une phrase précise, ce que cet exemple révèle sur la sensibilité de chacun de ces trois
indicateurs aux valeurs extrêmes.
\end{enumerate}

\end{problem}

\noindent Le désaccord entre moyenne et médiane n'est pas un artifice de manuel : c'est la
raison pour laquelle l'INSEE publie systématiquement les deux quand il décrit les salaires en
France, le salaire médian y étant nettement inférieur au salaire moyen. L'écart entre les deux
nombres est lui-même une information --- il mesure à quel point la distribution est tirée vers le
haut par une minorité. La moyenne appartient à la tradition des moindres carrés de Legendre et
Gauss ; la médiane fut défendue par Laplace, qui démontra qu'elle minimise l'erreur absolue. La
démonstration de cette propriété est STA-01, dans la bande Lycée.

\begin{refs}
\item Contexte : Moyenne arithmétique --- Wikipédia (\url{https://fr.wikipedia.org/wiki/Moyenne_arithm%C3%A9tique})
\item Contexte : Médiane (statistiques) --- Wikipédia (\url{https://fr.wikipedia.org/wiki/M%C3%A9diane_(statistiques)})
\item Contexte : Statistique descriptive --- Wikipédia (\url{https://fr.wikipedia.org/wiki/Statistique_descriptive})
\end{refs}

\bigskip

\begin{problem}[{L'axe qui ne part pas de zéro --- Collège}]
\textbf{STA-42.} \textbf{Problème.} Un journal local publie un diagramme en barres du nombre de cambriolages
enregistrés dans la ville sur cinq années consécutives : 1 010, 1 025, 1 040, 1 030 et 1 060.
L'axe vertical du graphique commence à 1 000 et s'arrête à 1 070. Le titre de l'article est \og{} La
délinquance flambe \fg{}.
\begin{enumerate}
\item[(a)] Calculez l'augmentation entre la première et la dernière année, en nombre de cambriolages puis
en pourcentage. Justifiez vos deux calculs et dites lequel des deux vous paraît le plus honnête à
publier.
\item[(b)] Tracez le même diagramme avec un axe vertical partant de 0. Décrivez précisément ce qui change
dans l'impression visuelle, et expliquez pourquoi les deux graphiques sont pourtant tous les deux
exacts.
\item[(c)] Un axe tronqué n'est pas toujours malhonnête. Donnez un exemple de série de données pour
laquelle partir de zéro serait au contraire une mauvaise idée, puis rédigez la règle que vous
proposeriez à ce journal pour distinguer les deux situations.
\end{enumerate}

\end{problem}

\noindent Darrell Huff a consacré à ce procédé un chapitre entier de \textit{How to Lie with
Statistics} (1954), l'un des livres de statistique les plus vendus de l'histoire ; il l'appelait
le graphique \og{} gee-whiz \fg{}, du nom de l'exclamation qu'il cherche à provoquer. Le point délicat, et
c'est celui de la question (c), est qu'aucune règle mécanique ne fonctionne : un axe partant de
zéro rendrait illisible une courbe de températures corporelles, et personne ne songerait à le
reprocher. Ce qui distingue l'honnêteté du mensonge n'est donc pas la position de l'axe mais
l'accord entre ce que le graphique fait voir et ce que le texte affirme.

\begin{refs}
\item Contexte : Visualisation de données --- Wikipédia (\url{https://fr.wikipedia.org/wiki/Visualisation_de_donn%C3%A9es})
\item Contexte : Diagramme à barres --- Wikipédia (\url{https://fr.wikipedia.org/wiki/Diagramme_%C3%A0_barres})
\item Contexte : Graphique trompeur --- Wikipédia (en anglais) (\url{https://en.wikipedia.org/wiki/Misleading_graph})
\end{refs}

\bigskip

\begin{problem}[{Le sondage de la cour de récréation --- Collège}]
\textbf{STA-43.} \textbf{Problème.} Camille veut connaître le sport préféré des 600 élèves de son collège. Un
mercredi à midi, elle interroge les 40 élèves présents sur le terrain de basket, dont 28 répondent
\og{} le basket \fg{}, et elle conclut : \og{} 70 \% des élèves du collège préfèrent le basket \fg{}.
\begin{enumerate}
\item[(a)] Énoncez au moins trois raisons distinctes pour lesquelles cette conclusion n'est pas fondée.
Pour chacune, dites exactement quelle étape du raisonnement de Camille elle attaque.
\item[(b)] Camille répond : \og{} j'ai interrogé 40 élèves, c'est beaucoup \fg{}. Expliquez pourquoi interroger
400 élèves au même endroit et au même moment ne corrigerait pas le défaut principal de son
enquête. Votre explication doit porter sur la façon dont les élèves interrogés ont été choisis.
\item[(c)] Concevez un protocole de sondage réellement applicable dans votre collège. Décrivez-le en cinq
lignes, justifiez chacun de vos choix, et terminez en indiquant ce que votre protocole ne garantit
toujours pas.
\end{enumerate}

\end{problem}

\noindent L'histoire des sondages est une suite d'échecs instructifs. En 1948, tous les
grands instituts américains annoncèrent la défaite de Harry Truman face à Thomas Dewey ; le
\textit{Chicago Tribune} imprima même sa une \og{} Dewey Defeats Truman \fg{} avant le dépouillement, et la
photographie de Truman brandissant le journal est restée célèbre. Les instituts avaient cessé de
sonder trop tôt et travaillaient sur des échantillons construits par quotas mal actualisés. Le cas
plus ancien et plus spectaculaire du \textit{Literary Digest}, qui recueillit deux millions de
réponses et se trompa quand même, est traité en STA-32 dans la bande Lycée : la leçon commune est
que la manière de choisir les personnes interrogées compte davantage que leur nombre.

\begin{refs}
\item Contexte : Sondage d'opinion --- Wikipédia (\url{https://fr.wikipedia.org/wiki/Sondage_d%27opinion})
\item Contexte : Biais de sélection --- Wikipédia (\url{https://fr.wikipedia.org/wiki/Biais_de_s%C3%A9lection})
\item Contexte : Dewey Defeats Truman --- Wikipédia (\url{https://fr.wikipedia.org/wiki/Dewey_Defeats_Truman})
\end{refs}

\bigskip

\begin{problem}[{Deux classes, une seule moyenne --- Collège}]
\textbf{STA-44.} \textbf{Problème.} Deux classes de neuf élèves ont passé le même contrôle, noté sur 20. Voici les
notes obtenues.

\noindent Classe A : 9, 10, 10, 11, 11, 11, 12, 12, 13 
Classe B : 3, 4, 5, 11, 13, 15, 16, 16, 16
\begin{enumerate}
\item[(a)] Vérifiez que les deux classes ont exactement la même moyenne, puis calculez pour chacune la
médiane et l'étendue. Justifiez chaque calcul.
\item[(b)] Le professeur doit dire laquelle des deux classes \og{} a mieux réussi \fg{}. Construisez d'abord
l'argument le plus solide en faveur de la classe A, puis l'argument le plus solide en faveur de la
classe B, chaque fois chiffres à l'appui. Concluez sur ce que la question posée a d'incomplet.
\item[(c)] Proposez un indicateur simple, calculable au collège, qui distinguerait ces deux séries mieux
que la moyenne. Justifiez qu'il les distingue effectivement, puis donnez un exemple de deux séries
que votre indicateur ne distingue pas non plus.
\end{enumerate}

\end{problem}

\noindent Une moyenne seule ne décrit jamais une série : c'est la première leçon de la
statistique descriptive, et elle a mis un siècle à se formaliser. On a longtemps mesuré la
dispersion par l'écart moyen ; c'est Karl Pearson qui, en 1893, a introduit le terme d'écart type
et imposé peu à peu la mesure que vous rencontrerez au lycée. Le choix entre plusieurs mesures de
dispersion n'a rien d'évident et fut sérieusement débattu --- la question (c) vous fait refaire ce
débat en miniature, y compris sa partie la plus honnête : reconnaître ce que votre indicateur
laisse encore passer.

\begin{refs}
\item Contexte : Série statistique --- Wikipédia (\url{https://fr.wikipedia.org/wiki/S%C3%A9rie_statistique})
\item Contexte : Médiane (statistiques) --- Wikipédia (\url{https://fr.wikipedia.org/wiki/M%C3%A9diane_(statistiques)})
\item Contexte : Karl Pearson --- Wikipédia (\url{https://fr.wikipedia.org/wiki/Karl_Pearson})
\end{refs}

\bigskip

\begin{problem}[{La moyenne des moyennes --- Collège}]
\textbf{STA-45.} \textbf{Problème.} Au dernier contrôle de mathématiques, la classe de 4e A, qui compte
28 élèves, a obtenu 12 de moyenne ; la classe de 4e B, qui en compte 12, a obtenu 16.
\begin{enumerate}
\item[(a)] Le principal annonce que \og{} la moyenne des deux classes réunies est 14 \fg{}. Montrez que c'est
faux, calculez la valeur correcte et justifiez votre méthode.
\item[(b)] Énoncez précisément la condition sous laquelle la moyenne des moyennes coïncide avec la
moyenne de l'ensemble. Démontrez-la dans le cas de deux groupes d'effectifs quelconques.
\item[(c)] Un élève a obtenu 8 de moyenne au premier trimestre, où il avait trois notes, et 16 au
deuxième, où il n'en avait qu'une. Il annonce 12 de moyenne sur les deux trimestres. Dites ce qui
est vrai et ce qui est faux dans son calcul, et déterminez en justifiant la valeur correcte.
\end{enumerate}

\end{problem}

\noindent Confondre la moyenne des moyennes et la moyenne globale est l'erreur
arithmétique la plus fréquemment publiée, et ce n'est pas une simple étourderie : c'est le
mécanisme même du paradoxe de Simpson, où deux groupes peuvent chacun avoir un meilleur taux de
réussite qu'un concurrent tout en ayant, réunis, un taux inférieur. La démonstration complète est
STA-02, dans la bande Lycée. Le système français des coefficients sur les bulletins scolaires est
d'ailleurs une moyenne pondérée déguisée : la question (c) montre que même sans coefficient
affiché, les effectifs jouent déjà ce rôle.

\begin{refs}
\item Contexte : Moyenne pondérée --- Wikipédia (\url{https://fr.wikipedia.org/wiki/Moyenne_pond%C3%A9r%C3%A9e})
\item Contexte : Paradoxe de Simpson --- Wikipédia (\url{https://fr.wikipedia.org/wiki/Paradoxe_de_Simpson})
\end{refs}

\bigskip

\begin{problem}[{Le diagramme circulaire du foyer --- Collège}]
\textbf{STA-46.} \textbf{Problème.} Le foyer socio-éducatif d'un collège publie un diagramme circulaire des
activités choisies par ses 240 inscrits. Les angles au centre des secteurs sont : théâtre 90$^\circ$,
échecs 60$^\circ$, journal 45$^\circ$, robotique 120$^\circ$, autres activités 45$^\circ$.
\begin{enumerate}
\item[(a)] Retrouvez l'effectif de chaque activité et justifiez la méthode de conversion que vous
employez pour passer d'un angle à un effectif. Vérifiez votre résultat.
\item[(b)] L'année suivante, le foyer compte 300 inscrits et le secteur \og{} robotique \fg{} occupe toujours
120$^\circ$. Le président en conclut que la robotique n'a pas progressé. Réfutez-le en distinguant
soigneusement effectif et fréquence, et dites ce que le diagramme permet malgré tout
d'affirmer.
\item[(c)] Le président souhaite mettre en valeur le journal du collège et propose de dessiner le
diagramme en perspective, le disque étant incliné. Expliquez géométriquement pourquoi cette
présentation fausse la lecture, et déterminez quels secteurs elle avantage.
\end{enumerate}

\end{problem}

\noindent Le diagramme circulaire est une invention de l'Écossais William Playfair, qui en
publia le premier exemple connu dans son \textit{Statistical Breviary} de 1801 ; on lui doit aussi le
diagramme en barres et la courbe chronologique, c'est-à-dire l'essentiel du vocabulaire graphique
que l'on emploie encore. Il a mauvaise réputation chez les spécialistes : Edward Tufte, dans
\textit{The Visual Display of Quantitative Information}, soutient que l'œil compare mal les angles et
que la seule chose pire qu'un diagramme circulaire est plusieurs diagrammes circulaires côte à
côte. La question (c) explique la moitié de cette mauvaise réputation.

\begin{refs}
\item Contexte : Diagramme circulaire --- Wikipédia (\url{https://fr.wikipedia.org/wiki/Diagramme_circulaire})
\item Contexte : William Playfair --- Wikipédia (\url{https://fr.wikipedia.org/wiki/William_Playfair})
\item Contexte : Edward Tufte --- Wikipédia (\url{https://fr.wikipedia.org/wiki/Edward_Tufte})
\end{refs}

\bigskip

\begin{problem}[{Effectif ou fréquence ? --- Collège, short}]
\textbf{STA-47.} \textbf{Problème.} Dans un collège de 800 élèves, 200 sont demi-pensionnaires ;
dans le collège voisin, 300 élèves sur 1 500 le sont. Un journaliste écrit : \og{} le second collège
compte plus de demi-pensionnaires, la demi-pension y est donc plus populaire \fg{} --- dites, en
justifiant par les effectifs puis par les fréquences, ce que cette phrase confond, et rédigez la
comparaison correcte.
\end{problem}

\noindent C'est pour cette raison que les statistiques officielles ne publient presque
jamais un effectif seul : l'INSEE et les agences sanitaires donnent des taux --- pour mille
habitants, pour cent mille --- précisément parce qu'un nombre brut mélange deux informations, la
grandeur du phénomène et la taille de la population où on le mesure. Un nombre de cas qui augmente
dans une population qui augmente ne dit rien tant qu'on n'a pas divisé. Séparer \og{} combien \fg{} de
\og{} quelle part \fg{} est le premier réflexe de lecture statistique, et l'un des plus souvent oubliés.

\begin{refs}
\item Contexte : Fréquence (statistiques) --- Wikipédia (\url{https://fr.wikipedia.org/wiki/Fr%C3%A9quence_(statistiques)})
\item Contexte : Pourcentage --- Wikipédia (\url{https://fr.wikipedia.org/wiki/Pourcentage})
\end{refs}

\bigskip

\begin{problem}[{Un pourcentage d'un pourcentage --- Collège, short}]
\textbf{STA-48.} \textbf{Problème.} Dans un collège, 40 \% des élèves sont en cycle 4 et, parmi
ces derniers, 30 \% apprennent l'allemand ; un délégué en conclut que \og{} 70 \% des élèves du collège
apprennent l'allemand en cycle 4 \fg{}. Débusquez l'erreur, déterminez en le justifiant le pourcentage
correct, et expliquez pourquoi la réponse ne dépend pas de l'effectif total du collège.
\end{problem}

\noindent Prendre un pourcentage d'un pourcentage, c'est multiplier des proportions, et
non les additionner --- une règle que l'arithmétique commerciale maîtrisait déjà bien avant la
statistique. Fibonacci l'applique aux intérêts et aux changes dans le \textit{Liber abaci} (1202), et
les écoles d'abaque de Toscane l'enseignaient aux futurs marchands du XIVe siècle. La
même erreur produit son cousin le plus célèbre : une hausse de 20 \% suivie d'une baisse de 20 \% ne
ramène pas au point de départ, et la question ci-dessus contient déjà la raison de ce
phénomène.

\begin{refs}
\item Contexte : Pourcentage --- Wikipédia (\url{https://fr.wikipedia.org/wiki/Pourcentage})
\item Contexte : Fréquence (statistiques) --- Wikipédia (\url{https://fr.wikipedia.org/wiki/Fr%C3%A9quence_(statistiques)})
\end{refs}

\bigskip

\begin{problem}[{La médiane d'un tableau d'effectifs --- Collège, short}]
\textbf{STA-49.} \textbf{Problème.} Dans une classe de 27 élèves, on a relevé le nombre de frères
et sœurs de chacun : 4 élèves en ont 0, 9 en ont 1, 7 en ont 2, 4 en ont 3, 2 en ont 4 et 1 en a 6.
Déterminez la médiane et la moyenne de cette série en justifiant votre méthode directement sur le
tableau d'effectifs, sans réécrire les 27 valeurs une à une, puis expliquez pourquoi la moyenne
n'est pas un nombre entier alors qu'aucun élève ne peut avoir un nombre non entier de frères et
sœurs.
\end{problem}

\noindent Le tableau d'effectifs est la forme sous laquelle les données arrivent
réellement --- un recensement ne vous remet pas une liste de soixante-sept millions de lignes, il
vous remet des comptages. Savoir lire une médiane et une moyenne directement dans un tableau
d'effectifs cumulés est donc une compétence de praticien avant d'être un exercice. Quant à la
dernière question, elle touche à quelque chose de profond : une moyenne n'est pas un individu de la
série, c'est un point d'équilibre, et rien n'oblige un point d'équilibre à exister parmi les
données.

\begin{refs}
\item Contexte : Médiane (statistiques) --- Wikipédia (\url{https://fr.wikipedia.org/wiki/M%C3%A9diane_(statistiques)})
\item Contexte : Série statistique --- Wikipédia (\url{https://fr.wikipedia.org/wiki/S%C3%A9rie_statistique})
\end{refs}

\bigskip

\begin{problem}[{Le pictogramme qui double --- Collège, short}]
\textbf{STA-50.} \textbf{Problème.} Pour illustrer que le budget du club de sport a doublé, un
graphiste dessine deux ballons dont le second a un diamètre deux fois plus grand que le premier.
Déterminez, calculs à l'appui, par quel facteur l'image exagère la hausse --- en distinguant le cas
où le lecteur compare des aires de disque et celui où il compare des volumes de ballon --- puis
proposez une construction correcte et justifiez-la.
\end{problem}

\noindent Darrell Huff consacre un chapitre de \textit{How to Lie with Statistics} (1954) à
ce procédé, qu'il tenait pour l'un des plus efficaces parce que le graphique reste, au sens strict,
véridique : le second dessin est bien deux fois plus grand que le premier \textit{dans une dimension}.
Le lecteur, lui, ne mesure pas des diamètres, il perçoit des surfaces ou des volumes, et l'écart
entre les deux est précisément ce que la question vous demande de chiffrer. C'est le seul exercice
de cette page où la géométrie et la statistique se rencontrent --- et ce n'est pas un hasard :
tromper avec un graphique, c'est presque toujours exploiter une différence entre ce qui est mesuré
et ce qui est vu.

\begin{refs}
\item Contexte : Pictogramme --- Wikipédia (\url{https://fr.wikipedia.org/wiki/Pictogramme})
\item Contexte : How to Lie with Statistics --- Wikipédia (en anglais) (\url{https://en.wikipedia.org/wiki/How_to_Lie_with_Statistics})
\item Contexte : Graphique trompeur --- Wikipédia (en anglais) (\url{https://en.wikipedia.org/wiki/Misleading_graph})
\end{refs}

\bigskip

\begin{problem}[{Combien d'élèves peuvent être au-dessus de la moyenne ? --- Collège}]
\textbf{STA-51.} \textbf{Problème.} On considère les notes d'un contrôle dans une classe de 25 élèves.
\begin{enumerate}
\item[(a)] Est-il possible que 24 des 25 élèves aient une note strictement supérieure à la moyenne de la
classe ? Construisez un exemple explicite, ou démontrez que c'est impossible.
\item[(b)] Démontrez qu'il est en revanche impossible que les 25 élèves soient tous strictement au-dessus
de la moyenne de la classe. Votre démonstration doit porter sur la somme des notes, et rester
valable pour un nombre quelconque d'élèves.
\item[(c)] Un journal titre : \og{} la moitié des collèges français sont en dessous de la moyenne nationale \fg{}.
Expliquez ce que ce titre a de vide, dites ce qu'il faudrait mesurer à la place pour que
l'information ait un contenu, et distinguez soigneusement le rôle de la moyenne et celui de la
médiane dans votre réponse.
\end{enumerate}

\end{problem}

\noindent L'idée que presque tout le monde peut être au-dessus de la moyenne porte un nom
en anglais, l'effet Lake Wobegon, d'après la ville imaginaire de l'humoriste Garrison Keillor \og{} où
tous les enfants sont au-dessus de la moyenne \fg{}. Elle n'est pas seulement drôle : le psychologue Ola
Svenson a publié en 1981 une enquête devenue classique où une large majorité d'automobilistes se
jugeaient plus prudents et plus habiles que la moyenne des conducteurs. La question (a) montre que
cela n'a rien d'impossible arithmétiquement, la question (b) dit où se trouve la vraie limite, et
la question (c) explique pourquoi confondre moyenne et médiane transforme une banalité
arithmétique en scandale de presse.

\begin{refs}
\item Contexte : Moyenne arithmétique --- Wikipédia (\url{https://fr.wikipedia.org/wiki/Moyenne_arithm%C3%A9tique})
\item Contexte : Supériorité illusoire --- Wikipédia (\url{https://fr.wikipedia.org/wiki/Sup%C3%A9riorit%C3%A9_illusoire})
\item Article : O. Svenson, \og{} Are we all less risky and more skillful than our fellow drivers? \fg{}, \textit{Acta Psychologica} 47 (1981), 143--148
\end{refs}

\bigskip

\begin{problem}[{Trois pourcentages dans la même case --- Collège}]
\textbf{STA-52.} \textbf{Problème.} On a interrogé 400 élèves d'un collège, dont 240 filles et 160 garçons, en
leur demandant s'ils pratiquent un sport en club. Parmi les filles, 156 répondent oui ; parmi les
garçons, 112.
\begin{enumerate}
\item[(a)] Recopiez et complétez le tableau à double entrée croisant le sexe et la pratique d'un sport en
club, marges comprises. Justifiez chaque case que vous ajoutez.
\item[(b)] À partir de la seule case \og{} filles pratiquant un sport en club \fg{}, on peut former les trois
quotients 156 sur 240, 156 sur 268 et 156 sur 400. Énoncez, pour chacun, la question précise à
laquelle il répond. Expliquez ensuite pourquoi confondre ces trois nombres est l'erreur la plus
courante dans la lecture d'un tel tableau.
\item[(c)] Le journal du collège titre : \og{} les filles font moins de sport que les garçons \fg{}. Dites ce que
les données permettent d'affirmer et ce qu'elles ne permettent pas, en distinguant explicitement
effectif et fréquence, puis nommez au moins une information manquante qui pourrait changer votre
jugement.
\end{enumerate}

\end{problem}

\noindent Le tableau à double entrée est le plus vieil instrument de la statistique
appliquée : John Graunt s'en sert déjà dans ses \textit{Natural and Political Observations Made upon the
Bills of Mortality} (1662), le premier travail de statistique démographique, pour croiser causes
de décès et années. Sa simplicité est trompeuse --- la confusion entre les trois quotients de la
question (b) est la racine commune de la plupart des erreurs d'interprétation, depuis les tests
médicaux mal lus jusqu'au paradoxe de Simpson (STA-02, bande Lycée). Apprendre à dire à voix haute
\og{} sur quoi je divise, et pourquoi \fg{} règle une bonne partie de ces malentendus.

\begin{refs}
\item Contexte : Tableau de contingence --- Wikipédia (\url{https://fr.wikipedia.org/wiki/Tableau_de_contingence})
\item Contexte : John Graunt --- Wikipédia (\url{https://fr.wikipedia.org/wiki/John_Graunt})
\item Contexte : Fréquence (statistiques) --- Wikipédia (\url{https://fr.wikipedia.org/wiki/Fr%C3%A9quence_(statistiques)})
\end{refs}

\bigskip


\section*{Lycée}

\begin{problem}[{Moyenne contre médiane --- Lycée}]
\textbf{STA-01.} \textbf{Problème.} Pour des réels $ x_1, \dots, x_n $, on pose
\[ f(c) = \sum_{i=1}^{n} (x_i - c)^2 , \qquad g(c) = \sum_{i=1}^{n} |x_i - c| . \]
\begin{enumerate}
\item[(a)] Démontrez que $ f $ atteint son minimum en un unique point, la moyenne des nombres.
\item[(b)] Démontrez que $ g $ atteint son minimum en toute médiane (tout point ayant au moins la
moitié des données de chaque côté), et déterminez exactement quand le minimiseur est unique.
\item[(c)] Construisez une liste de dix \og{} revenus des ménages \fg{} dont la moyenne dépasse dix fois la
médiane, et argumentez laquelle des deux est le \og{} revenu typique \fg{} honnête pour votre liste ---
avec une raison mathématique, pas un slogan.
\end{enumerate}

\end{problem}

\noindent La tension entre les deux moyennes est ancienne : la moyenne appartient à la
tradition des moindres carrés de Legendre et Gauss, tandis que la médiane fut défendue par Laplace
(qui démontra qu'elle minimise l'erreur absolue) puis par Edgeworth. Le choix tranche encore
aujourd'hui des débats publics --- revenus, prix de l'immobilier, délais d'attente à l'hôpital ---
parce que la moyenne court après les valeurs extrêmes et la médiane non. La partie (b) est une
véritable démonstration sur une fonction affine par morceaux plutôt que dérivable, ce qui est
exactement pourquoi elle mérite d'être rédigée.

\begin{refs}
\item Contexte : Médiane (statistiques) --- Wikipédia (\url{https://fr.wikipedia.org/wiki/M%C3%A9diane_(statistiques)})
\item Lecture : D. Freedman, R. Pisani \& R. Purves, \textit{Statistics}, ch. 4
\item Cours : MIT OCW 18.650 Statistics for Applications (\url{https://ocw.mit.edu/courses/18-650-statistics-for-applications-fall-2016/})
\end{refs}

\bigskip

\begin{problem}[{Construisez votre propre paradoxe de Simpson --- Lycée}]
\textbf{STA-02.} \textbf{Problème.} Deux hôpitaux, A et B, traitent chacun exactement 1000 patients, chaque
patient étant classé comme cas bénin ou cas grave.
\begin{enumerate}
\item[(a)] Construisez des données entières telles que : l'hôpital A ait un taux de survie strictement
supérieur à celui de B chez les cas bénins, A ait un taux de survie strictement supérieur à celui
de B chez les cas graves, et pourtant B ait un taux de survie global strictement supérieur.
\item[(b)] Démontrez le résultat d'impossibilité suivant : si les deux hôpitaux traitent le \textit{même
nombre} de cas bénins (et donc le même nombre de cas graves), aucune inversion de ce genre ne
peut se produire.
\end{enumerate}

\end{problem}

\noindent L'inversion porte le nom d'E. H. Simpson (1951), bien que Karl Pearson et Udny
Yule l'aient remarquée un demi-siècle plus tôt. Le cas réel le plus célèbre est celui des données
d'admission en troisième cycle à Berkeley en 1973, analysées par Bickel, Hammel et O'Connell :
l'université semblait globalement défavoriser les femmes, alors que département par département ce
n'était pas le cas --- les femmes avaient postulé aux départements les plus sélectifs. Le paradoxe
est un théorème sur les moyennes pondérées, et votre démonstration d'impossibilité montrera
exactement d'où vient la liberté d'inverser.

\begin{refs}
\item Contexte : Paradoxe de Simpson --- Wikipédia (\url{https://fr.wikipedia.org/wiki/Paradoxe_de_Simpson})
\item Lecture : J. Pearl \& D. Mackenzie, \textit{The Book of Why}, ch. 6
\item Article : P. J. Bickel, E. A. Hammel \& J. W. O'Connell, \og{} Sex bias in graduate admissions: Data from Berkeley \fg{}, \textit{Science} 187 (1975), 398--404
\end{refs}

\bigskip

\begin{problem}[{Corrélation n'est pas causalité --- sous forme de théorème --- Lycée}]
\textbf{STA-03.} \textbf{Problème.} La corrélation empirique $ r $ de données appariées mesure l'association
linéaire et vérifie toujours $ -1 \le r \le 1 $ (vous pouvez l'admettre).
\begin{enumerate}
\item[(a)] Soient $ Z, U, V $ des quantités aléatoires indépendantes, chacune d'espérance 0 et de
variance 1, et posons $ X = Z + U $ et $ Y = Z + V $. Montrez que la corrélation
(théorique) de $ X $ et $ Y $ vaut $ 1/2 $, alors que ni $ X $ ni $ Y $ n'influence
l'autre --- ils partagent seulement la cause commune $ Z $.
\item[(b)] Construisez une variable aléatoire $ X $ prenant un nombre fini de valeurs, telle que
$ X $ et $ Y = X^2 $ aient une corrélation exactement nulle bien que $ Y $ soit entièrement
déterminée par $ X $. Démontrez les deux affirmations.
\item[(c)] Concluez, en un paragraphe soigneux, sur ce qu'une corrélation non nulle établit et n'établit
pas, et sur ce qu'une corrélation nulle établit et n'établit pas.
\end{enumerate}

\end{problem}

\noindent \og{} Corrélation n'implique pas causalité \fg{} est d'ordinaire laissé à l'état de
proverbe ; il devient ici deux petits théorèmes. La partie (a) est le \textit{facteur de confusion} :
les ventes de glaces et les noyades sont corrélées parce que l'été entraîne les deux. La partie (b)
montre le piège réciproque : corrélation nulle n'est pas indépendance. Le coefficient de
corrélation lui-même est dû à Francis Galton et Karl Pearson dans les années 1880--1890 ; la
mathématique moderne des conditions sous lesquelles une corrélation \textit{peut} soutenir une
affirmation causale est l'objet de l'inférence causale, l'un des domaines les plus vivants de la
statistique d'aujourd'hui.

\begin{refs}
\item Contexte : Corrélation n'implique pas causalité --- Wikipédia (\url{https://fr.wikipedia.org/wiki/Corr%C3%A9lation_n'implique_pas_causalit%C3%A9})
\item Contexte : Coefficient de corrélation de Pearson --- Wikipédia (en anglais) (\url{https://en.wikipedia.org/wiki/Pearson_correlation_coefficient})
\item Lecture : D. Freedman, R. Pisani \& R. Purves, \textit{Statistics}, ch. 9--12
\end{refs}

\bigskip

\begin{problem}[{Pourquoi les fils des pères grands sont plus petits --- Lycée, short}]
\textbf{STA-19.} \textbf{Problème.} Soient $ X $ et $ Y $ les tailles standardisées du père
et du fils (chacune d'espérance 0 et d'écart type 1), de corrélation $ r $, et supposons la
moyenne conditionnelle linéaire : $ \mathbb{E}[Y \mid X = x] = a + bx $. Démontrez que
$ a = 0 $ et $ b = r $, de sorte que $ |\mathbb{E}[Y \mid X = x]| <|x| $ dès que
$ |r| <1 $ et $ x \ne 0 $, puis expliquez --- par un argument de symétrie, pas par un slogan ---
pourquoi cela ne signifie \textit{pas} que la dispersion des tailles se réduit de génération en
génération.
\end{problem}

\noindent Francis Galton découvrit l'effet dans les années 1880 et l'appela \og{} régression
vers la médiocrité \fg{} ; toute la discipline moderne de l'analyse de régression tire son nom de cette
seule observation biologique. Le point de la seconde moitié du problème est que le même argument
mené à l'envers prédit que les fils grands avaient des pères plus petits, ce qu'aucune théorie de
l'hérédité ne peut expliquer --- l'effet est une propriété de la corrélation imparfaite, pas de la
biologie. C'est aussi le piège classique de l'évaluation des interventions : les écoles les moins
performantes, les patients les plus malades et les pilotes les plus malchanceux s'améliorent tout
seuls, et quelque chose d'autre en récolte toujours le crédit.

\begin{refs}
\item Contexte : Régression vers la moyenne --- Wikipédia (\url{https://fr.wikipedia.org/wiki/R%C3%A9gression_vers_la_moyenne})
\item Article : F. Galton, \og{} Regression towards mediocrity in hereditary stature \fg{}, \textit{Journal of the Anthropological Institute} 15 (1886), 246--263
\item Lecture : D. Freedman, R. Pisani \& R. Purves, \textit{Statistics}, ch. 10
\end{refs}

\bigskip

\begin{problem}[{Une garantie qui n'a besoin d'aucune courbe en cloche --- Lycée, short}]
\textbf{STA-20.} \textbf{Problème.} Soit $ x_1, \dots, x_n $ une liste quelconque de nombres,
de moyenne $ \bar{x} $ et d'écart type
$ s = \sqrt{\frac{1}{n}\sum_i (x_i - \bar{x})^2} $, et soit $ k >1 $. Démontrez que la
proportion de la liste située à une distance au moins $ ks $ de $ \bar{x} $ est au plus
$ 1/k^2 $, et montrez que la borne ne peut pas être améliorée en exhibant une liste de huit
nombres pour laquelle cette proportion vaut exactement $ 1/4 $ lorsque $ k = 2 $.
\end{problem}

\noindent C'est l'inégalité de Bienaymé-Tchebychev dépouillée de toute probabilité : la
démonstration tient en une ligne d'arithmétique --- chaque terme de la somme des carrés est positif,
alors jetez les petits. La fameuse règle \og{} 68-95-99,7 \fg{} est bien plus fine, mais elle ne vaut que
pour la courbe normale ; la borne de Tchebychev vaut pour des données de revenus, des magnitudes de
séismes et n'importe quelle liste que vous rencontrerez, ce qui explique que Tchebychev s'en soit
servi en 1867 pour démontrer la loi des grands nombres sans rien supposer sur la forme de la
loi.

\begin{refs}
\item Contexte : Inégalité de Bienaymé-Tchebychev --- Wikipédia (\url{https://fr.wikipedia.org/wiki/In%C3%A9galit%C3%A9_de_Bienaym%C3%A9-Tchebychev})
\item Contexte : Règle 68-95-99,7 --- Wikipédia (\url{https://fr.wikipedia.org/wiki/R%C3%A8gle_68-95-99,7})
\item Lecture : D. Freedman, R. Pisani \& R. Purves, \textit{Statistics}, ch. 4--5
\end{refs}

\bigskip

\begin{problem}[{Compter des poissons que l'on ne peut pas compter --- Lycée}]
\textbf{STA-21.} \textbf{Problème.} Un lac contient un nombre inconnu $ N $ de poissons. Vous en capturez
$ K $, vous les marquez et vous les relâchez ; plus tard vous capturez un second échantillon de
$ n $ poissons, dont $ k $ se révèlent marqués. On suppose que tout sous-ensemble de $ n $
poissons a la même probabilité d'être le second échantillon.
\begin{enumerate}
\item[(a)] Écrivez la probabilité d'observer exactement $ k $ poissons marqués (une probabilité
hypergéométrique), en fonction de l'inconnue $ N $.
\item[(b)] En examinant le rapport $ L(N)/L(N-1) $ de vraisemblances consécutives, démontrez que cette
probabilité est maximale en $ \hat{N} = \lfloor Kn/k \rfloor $ lorsque $ k \ge 1 $ ---
l'estimateur du maximum de vraisemblance est exactement la règle de proportionnalité naïve
$ k/n \approx K/N $.
\item[(c)] Montrez que $ \hat{N} $ est inutilisable tel quel : $ P(k = 0) >0 $ pour tout
$ N \ge K + n $, et sur cet événement l'estimation est infinie. Démontrez que l'estimateur de
Chapman
\[ \tilde{N} = \frac{(K+1)(n+1)}{k+1} - 1 \]
est exactement sans biais dès que $ K + n \ge N $. (Une identité de type Vandermonde pour les
coefficients binomiaux fait le travail.)
\item[(d)] Calculez les deux estimations pour $ K = 100 $, $ n = 100 $, $ k = 20 $. Énumérez ensuite
les hypothèses de modélisation --- population fermée, aucune marque perdue, tous les poissons
également capturables --- et pour chacune dites dans quel sens son échec ferait dévier
l'estimation.
\end{enumerate}

\end{problem}

\noindent Laplace utilisa la même arithmétique en 1802 pour estimer la population de la
France à partir des registres de naissances, et la méthode fut réinventée pour les animaux par
C. G. J. Petersen pour les poissons (1896) et F. C. Lincoln pour les oiseaux d'eau bagués (1930) ;
la correction de biais de (c) est celle de D. G. Chapman, en 1951. La technique sert aujourd'hui
bien au-delà de l'écologie : estimer le sous-dénombrement d'un recensement, le nombre de défauts
logiciels non découverts et --- en croisant plusieurs listes incomplètes --- le nombre de victimes d'un
conflit armé. C'est l'exemple le plus net, en statistique élémentaire, de la façon d'apprendre la
taille d'une chose que l'on ne pourra jamais voir en entier.

\begin{refs}
\item Contexte : Capture-marquage-recapture --- Wikipédia (\url{https://fr.wikipedia.org/wiki/Capture-marquage-recapture})
\item Contexte : Loi hypergéométrique --- Wikipédia (\url{https://fr.wikipedia.org/wiki/Loi_hyperg%C3%A9om%C3%A9trique})
\item Lecture : G. A. F. Seber, \textit{The Estimation of Animal Abundance and Related Parameters}, 2e éd.
\end{refs}

\bigskip

\begin{problem}[{La dame qui déguste du thé --- Lycée, short}]
\textbf{STA-31.} \textbf{Problème.} On prépare huit tasses de thé, quatre où le lait a été versé
en premier et quatre où c'est le thé ; on informe la dégustatrice qu'il y en a quatre de chaque
sorte et elle doit désigner les quatre tasses au lait versé en premier, de sorte que sous
l'hypothèse nulle où elle devine au hasard, les $ \binom{8}{4} = 70 $ sélections possibles sont
équiprobables. Démontrez que la probabilité de désigner les quatre correctement par hasard vaut
$ 1/70 $, calculez la probabilité d'en trouver au moins trois sur quatre, et déduisez-en qu'au
seuil de 5 \% rien de moins qu'un sans-faute ne peut compter comme preuve --- puis dites ce qui change
si l'expérience est menée avec six tasses, trois de chaque sorte.
\end{problem}

\noindent Ronald Fisher ouvre \textit{The Design of Experiments} (1935) par cette
expérience, et elle a réellement eu lieu : à Rothamsted, dans les années 1920, l'algologue Muriel
Bristol affirmait pouvoir sentir au goût si le lait avait été versé dans la tasse en premier, et
Fisher conçut le test sur-le-champ. Ce chapitre est le premier endroit où la randomisation,
l'hypothèse nulle et la valeur p exacte apparaissent ensemble, et le dénombrement que vous venez de
faire est le \textit{test exact de Fisher} pour un tableau $ 2 \times 2 $ --- pas d'approximation
normale, pas de théorie asymptotique, rien que de la combinatoire. Remarquez ce que révèle la
dernière partie : le plus petit effet détectable est fixé par le plan d'expérience avant qu'une
seule tasse ne soit versée, et c'est pourquoi Fisher insistait pour que l'on consulte le
statisticien avant l'expérience plutôt qu'après.

\begin{refs}
\item Contexte : La dame dégustant du thé --- Wikipédia (\url{https://fr.wikipedia.org/wiki/The_lady_tasting_tea})
\item Contexte : Test exact de Fisher --- Wikipédia (\url{https://fr.wikipedia.org/wiki/Test_exact_de_Fisher})
\item Lecture : R. A. Fisher, \textit{The Design of Experiments} (1935), ch. 2
\item Lecture : D. Salsburg, \textit{The Lady Tasting Tea: How Statistics Revolutionized Science in the Twentieth Century} (2001)
\end{refs}

\bigskip

\begin{problem}[{Deux millions de mauvaises réponses --- Lycée}]
\textbf{STA-32.} \textbf{Problème.} Une population comporte une fraction $ p $ d'électeurs favorables au candidat
R et $ 1 - p $ favorables au candidat L. Un magazine envoie des bulletins par la poste à une
immense liste de noms. Indépendamment de tout le reste, un électeur favorable à R renvoie le
bulletin avec probabilité $ r_1 $, et un électeur favorable à L le renvoie avec probabilité
$ r_0 $ ; le magazine publie la fraction de partisans de R parmi les bulletins qui lui
reviennent.
\begin{enumerate}
\item[(a)] Démontrez que lorsque le nombre de bulletins envoyés croît, la fraction publiée se stabilise
non pas en $ p $ mais en
\[ p^{*} = \frac{p\, r_1}{p\, r_1 + (1-p)\, r_0} , \]
et que $ p^{*} = p $ si et seulement si $ r_1 = r_0 $. Remarquez que le nombre de bulletins
n'apparaît nulle part dans $ p^{*} $ : l'erreur ne diminue pas quand le sondage grossit.
\item[(b)] Injectez les chiffres historiques. En 1936, le \textit{Literary Digest} envoya environ dix
millions de bulletins, en reçut environ 2,38 millions en retour, et prédit que Landon obtiendrait
57 \% des voix ; Roosevelt en obtint 60,8 \%. À l'aide de (a), déterminez le rapport
$ r_1/r_0 $ des taux de réponse qui expliquerait à lui seul l'écart, et commentez la modestie de
ce différentiel.
\item[(c)] Montrez que l'erreur d'échantillonnage n'était pas le problème. Pour un échantillon réellement
aléatoire de taille $ n $ tiré de la population, démontrez que l'écart type de la fraction
observée vaut $ \sqrt{p(1-p)/n} \le 1/(2\sqrt{n}) $, et évaluez-le pour
$ n = 2\,380\,000 $. Comparez le résultat à l'erreur de 20 points du magazine, puis trouvez la
taille d'échantillon à laquelle la seule erreur d'échantillonnage aléatoire atteint un point de
pourcentage.
\item[(d)] Le \textit{Digest} a commis deux erreurs distinctes : il a bâti son fichier d'adresses à partir
d'annuaires téléphoniques, de listes de clubs et d'immatriculations d'automobiles, et il a accepté
la réponse de quiconque voulait bien répondre. Expliquez, à l'aide de (a), pourquoi chacune est
fatale à elle seule, et démontrez qu'envoyer vingt millions de bulletins au lieu de dix n'aurait
remédié ni à l'une ni à l'autre.
\end{enumerate}

\end{problem}

\noindent Le \textit{Literary Digest} avait correctement annoncé toutes les élections
présidentielles depuis 1920 et était si confiant dans le volume brut qu'il publia sa prévision de
1936 comme une quasi-certitude ; George Gallup, travaillant sur un échantillon d'environ cinquante
mille personnes choisies pour représenter l'électorat, annonça le résultat correct, et le magazine
disparut en moins de deux ans. La reconstitution de Peverill Squire en 1988 a identifié la
non-réponse comme le plus grave des deux échecs : les partisans de Roosevelt renvoyaient tout
simplement moins de bulletins. La leçon n'est pas seulement historique. Xiao-Li Meng a montré en
2018 que l'erreur d'une moyenne d'échantillon se factorise en un terme de qualité des données, un
terme de quantité de données et la variance de la population, de sorte qu'une corrélation minuscule
entre le fait de répondre et la préférence rend un échantillon énorme aussi mauvais qu'un petit
échantillon aléatoire --- son \og{} paradoxe des mégadonnées \fg{}, appliqué à l'élection américaine de 2016,
est STA-32(a) réécrit pour le XXIe siècle.

\begin{refs}
\item Contexte : The Literary Digest --- Wikipédia (en anglais) (\url{https://en.wikipedia.org/wiki/The_Literary_Digest})
\item Contexte : Échantillon biaisé --- Wikipédia (\url{https://fr.wikipedia.org/wiki/%C3%89chantillon_biais%C3%A9})
\item Article : P. Squire, \og{} Why the 1936 Literary Digest poll failed \fg{}, \textit{Public Opinion Quarterly} 52 (1988), 125--133
\item Article : X.-L. Meng, \og{} Statistical paradises and paradoxes in big data (I) \fg{}, \textit{Annals of Applied Statistics} 12 (2018), 685--726
\end{refs}

\bigskip


\section*{Licence}

\begin{problem}[{Les moindres carrés et les équations normales --- Licence}]
\textbf{STA-04.} \textbf{Problème.} Étant donné des points $ (x_1, y_1), \dots, (x_n, y_n) $ dont les $ x_i $
ne sont pas tous égaux, on ajuste une droite $ y = a + bx $ en minimisant
\[ S(a,b) = \sum_{i=1}^{n} (y_i - a - b x_i)^2 . \]
\begin{enumerate}
\item[(a)] Établissez les équations normales en annulant les deux dérivées partielles
$ \partial S/\partial a $ et $ \partial S/\partial b $, résolvez-les, et démontrez que le point
critique est l'unique minimum global.
\item[(b)] Montrez que la droite ajustée passe toujours par le point moyen
$ (\bar{x}, \bar{y}) $.
\item[(c)] Sous forme matricielle, $ X $ désignant la matrice de plan $ n \times p $ de rang maximal
en colonnes, montrez que minimiser $ \lVert y - X\beta \rVert^2 $ conduit aux équations normales
$ X^{\top} X \beta = X^{\top} y $. Interprétez la solution géométriquement --- le vecteur ajusté
$ X\hat{\beta} $ est la projection orthogonale de $ y $ sur l'espace des colonnes de $ X $ ---
et démontrez cette interprétation.
\end{enumerate}

\end{problem}

\noindent Les moindres carrés sont l'algorithme fondateur de la statistique. Legendre les
publia en 1805 dans un appendice sur les orbites cométaires ; Gauss affirma les employer depuis
1795 et les relia en 1809 à la loi normale --- la querelle de priorité fut âpre. Le premier triomphe
de la méthode fut de retrouver l'astéroïde Cérès perdu à partir d'une poignée d'observations. Les
équations normales que vous établissez ici sont résolues des millions de fois par seconde au cœur
de tous les systèmes de régression, de prévision et d'apprentissage automatique qui existent.

\begin{refs}
\item Contexte : Méthode des moindres carrés --- Wikipédia (\url{https://fr.wikipedia.org/wiki/M%C3%A9thode_des_moindres_carr%C3%A9s})
\item Lecture : G. Casella \& R. L. Berger, \textit{Statistical Inference}, 2e éd., ch. 11--12
\item Cours : MIT OCW 18.650 Statistics for Applications (\url{https://ocw.mit.edu/courses/18-650-statistics-for-applications-fall-2016/})
\end{refs}

\bigskip

\begin{problem}[{Le maximum de vraisemblance, quatre fois --- Licence}]
\textbf{STA-05.} \textbf{Problème.} Des observations indépendantes $ x_1, \dots, x_n $ sont tirées de chacun des
modèles ci-dessous. Pour chacun des cas (a)--(d), écrivez la vraisemblance, établissez l'estimateur
du maximum de vraisemblance (EMV), et vérifiez que c'est bien un maximum.
\begin{enumerate}
\item[(a)] Bernoulli de probabilité de succès $ p $.
\item[(b)] Poisson d'espérance $ \lambda $.
\item[(c)] Exponentielle de taux $ \lambda $.
\item[(d)] Normale $ N(\mu, \sigma^2) $ avec les deux paramètres inconnus.
\item[(e)] Déterminez, pour chaque estimateur trouvé, s'il est sans biais ; là où il est biaisé, calculez
le biais exactement.
\item[(f)] Démontrez la propriété d'invariance : si $ \hat{\theta} $ est un EMV de $ \theta $ et si
$ g $ est une bijection, alors $ g(\hat{\theta}) $ est un EMV de $ g(\theta) $.
\end{enumerate}

\end{problem}

\noindent Le maximum de vraisemblance est la création de R. A. Fisher --- esquissé dans un
article d'étudiant de 1912 et développé dans son grand mémoire de 1922, qui introduisit aussi les
mots \og{} paramètre \fg{}, \og{} statistique \fg{}, \og{} exhaustivité \fg{} et \og{} efficacité \fg{} dans leur sens moderne. Le
principe --- rendre les données aussi probables que possible --- a l'air d'une heuristique, et pourtant
il produit des estimateurs asymptotiquement optimaux sous des conditions très générales (voir
STA-14). La partie (e) réserve une petite surprise : l'un des estimateurs les plus naturels de la
statistique est biaisé, ce qui mène directement au problème suivant.

\begin{refs}
\item Contexte : Maximum de vraisemblance --- Wikipédia (\url{https://fr.wikipedia.org/wiki/Maximum_de_vraisemblance})
\item Lecture : G. Casella \& R. L. Berger, \textit{Statistical Inference}, 2e éd., ch. 7
\item Article : R. A. Fisher, \og{} On the mathematical foundations of theoretical statistics \fg{}, \textit{Philosophical Transactions of the Royal Society A} 222 (1922), 309--368
\end{refs}

\bigskip

\begin{problem}[{Pourquoi diviser par n $-$ 1 ? --- Licence}]
\textbf{STA-06.} \textbf{Problème.} Soient $ x_1, \dots, x_n $ des tirages indépendants d'une loi quelconque
d'espérance $ \mu $ et de variance finie $ \sigma^2 $, et soit $ \bar{x} $ la moyenne
empirique. On définit la variance empirique
\[ S^2 = \frac{1}{n-1} \sum_{i=1}^{n} (x_i - \bar{x})^2 . \]
\begin{enumerate}
\item[(a)] Démontrez l'identité
\[ \sum_{i=1}^{n} (x_i - \bar{x})^2 = \sum_{i=1}^{n} (x_i - \mu)^2 - n(\bar{x} - \mu)^2 , \]
et servez-vous-en pour démontrer $ \mathbb{E}[S^2] = \sigma^2 $.
\item[(b)] Montrez que diviser par $ n $ donne au contraire l'espérance $ \sigma^2 (n-1)/n $, de sorte
que l'estimateur \og{} naturel \fg{} sous-estime systématiquement la variance. Expliquez avec vos mots
pourquoi mesurer la dispersion autour de $ \bar{x} $ plutôt qu'autour de $ \mu $ fait perdre
exactement l'information d'une observation.
\item[(c)] Montrez que l'absence de biais ne survit pas à la racine carrée : $ S $ n'est \textit{pas} un
estimateur sans biais de $ \sigma $ (supposez que la loi n'est pas concentrée en un point).
\end{enumerate}

\end{problem}

\noindent Le n $-$ 1 s'appelle la correction de Bessel, d'après l'astronome Friedrich
Bessel, et c'est la première rencontre honnête de la plupart des étudiants avec la notion de biais.
L'idée plus profonde --- que les résidus $ x_i - \bar{x} $ vivent dans un espace de dimension
$ (n-1) $, à l'origine de l'expression \og{} degrés de liberté \fg{} --- fut rendue précise par Fisher dans
les années 1920. La partie (c) est un avertissement utile : l'absence de biais est une propriété
d'un paramétrage donné, pas une loi de la nature.

\begin{refs}
\item Contexte : Correction de Bessel --- Wikipédia (\url{https://fr.wikipedia.org/wiki/Correction_de_Bessel})
\item Lecture : G. Casella \& R. L. Berger, \textit{Statistical Inference}, 2e éd., ch. 5
\end{refs}

\bigskip

\begin{problem}[{Ce que promet un intervalle de confiance --- Licence}]
\textbf{STA-07.} \textbf{Problème.} Soient $ x_1, \dots, x_n $ des observations indépendantes
$ N(\mu, \sigma^2) $ avec $ \sigma $ connu.
\begin{enumerate}
\item[(a)] Démontrez l'énoncé de couverture : l'intervalle aléatoire
$ \bar{x} \pm 1{,}96\,\sigma/\sqrt{n} $ contient $ \mu $ avec probabilité 0,95 --- la
probabilité portant sur le tirage des données, avant que celles-ci ne soient vues, et cela quelle
que soit la vraie valeur de $ \mu $.
\item[(b)] Expliquez précisément pourquoi la phrase \og{} ayant calculé l'intervalle $ [2{,}1 ;\,3{,}4] $,
le paramètre $ \mu $ s'y trouve avec probabilité 95 \% \fg{} n'est pas autorisée par (a), et quelle
est la bonne interprétation fréquentiste sur des expériences répétées.
\item[(c)] Voici maintenant l'exemple d'école qui met en garde. Un paramètre $ \theta $ est inconnu ;
deux observations indépendantes $ X_1, X_2 $ valent chacune $ \theta - 1 $ ou
$ \theta + 1 $, chacune avec probabilité $ 1/2 $. On définit l'ensemble de confiance
$ C $ : si $ X_1 \ne X_2 $, on pose $ C = \{ (X_1 + X_2)/2 \} $ ; si $ X_1 = X_2 $, on pose
$ C = \{ X_1 - 1 \} $. Démontrez que $ C $ couvre $ \theta $ avec probabilité exactement
75 \%, et pourtant qu'après avoir vu les données vous êtes soit logiquement certain que
$ C = \{\theta\} $, soit assuré que la couverture n'est qu'un coup de pile ou face. Que dit cela
sur le fait d'attacher le niveau de confiance pré-expérimental à un intervalle particulier une fois
réalisé ?
\end{enumerate}

\end{problem}

\noindent Les intervalles de confiance ont été introduits par Jerzy Neyman dans un article
de 1937 qui mérite encore d'être lu pour sa rigueur philosophique : le niveau de confiance décrit
la performance à long terme de la \textit{procédure}, pas la probabilité d'un intervalle réalisé
particulier. L'exemple (c) --- un exemple d'enseignement classique, discuté par Berger, Wolpert,
Wasserman et d'autres --- montre que les deux choses peuvent diverger spectaculairement. Les énoncés
sur la probabilité d'un paramètre sachant les données relèvent de l'inférence bayésienne (STA-10),
qui fait d'autres hypothèses pour se les autoriser.

\begin{refs}
\item Contexte : Intervalle de confiance --- Wikipédia (\url{https://fr.wikipedia.org/wiki/Intervalle_de_confiance})
\item Lecture : L. Wasserman, \textit{All of Statistics}, ch. 6
\item Article : J. Neyman, \og{} Outline of a theory of statistical estimation based on the classical theory of probability \fg{}, \textit{Philosophical Transactions of the Royal Society A} 236 (1937), 333--380
\end{refs}

\bigskip

\begin{problem}[{La logique des valeurs p --- Licence}]
\textbf{STA-08.} \textbf{Problème.} Soit $ T $ une statistique de test de loi continue sous l'hypothèse nulle
$ H_0 $, les grandes valeurs comptant contre $ H_0 $, et soit la valeur p
$ p = P_0(T \ge t_{\mathrm{obs}}) $.
\begin{enumerate}
\item[(a)] Démontrez que sous $ H_0 $ la valeur p suit la loi uniforme sur $ (0,1) $, et déduisez-en
que la règle \og{} rejeter lorsque $ p \le \alpha $ \fg{} a une probabilité d'erreur de première espèce
exactement égale à $ \alpha $.
\item[(b)] Expliquez soigneusement pourquoi $ p $ n'est pas la probabilité que $ H_0 $ soit vraie, et
pourquoi $ 1 - p $ n'est pas la probabilité que l'alternative soit vraie.
\item[(c)] Construisez un exemple explicite où la preuve empirique paraît accablante et où pourtant
l'hypothèse nulle survit à un audit bayésien : $ H_0\colon \theta = 0 $ contre
$ H_1\colon \theta \ne 0 $ pour l'espérance de données $ N(\theta, 1) $, avec une probabilité a
priori $ 1/2 $ sur $ H_0 $ et, sous $ H_1 $, une loi a priori normale sur $ \theta $.
Montrez que pour $ n $ grand on peut avoir $ p = 0{,}05 $ alors que la probabilité a posteriori
de $ H_0 $ dépasse $ 1/2 $. (C'est le paradoxe de Lindley ; le problème consiste à le poser
précisément.)
\end{enumerate}

\end{problem}

\noindent Les tests de significativité descendent de Fisher, qui voyait la valeur p comme
une mesure grossière de la surprise, et de Neyman et Pearson, qui ont refondu le test en une
décision à taux d'erreur contrôlés ; le rituel moderne du p $\le$ 0,05 est une fusion inconfortable des
deux qu'aucun n'approuverait. Les mauvaises lectures de (b) sont si répandues que l'American
Statistical Association a publié une déclaration officielle à leur sujet en 2016. Le paradoxe de
Lindley (c), publié en 1957, montre que \og{} significatif à 5 \% \fg{} et \og{} probablement vrai \fg{} peuvent
pointer dans des directions opposées.

\begin{refs}
\item Contexte : Valeur p --- Wikipédia (\url{https://fr.wikipedia.org/wiki/Valeur_p})
\item Contexte : Paradoxe de Lindley --- Wikipédia (en anglais) (\url{https://en.wikipedia.org/wiki/Lindley%27s_paradox})
\item Lecture : L. Wasserman, \textit{All of Statistics}, ch. 10
\item Article : R. L. Wasserstein \& N. A. Lazar, \og{} The ASA's statement on p-values \fg{}, \textit{The American Statistician} 70 (2016), 129--133
\end{refs}

\bigskip

\begin{problem}[{Le lemme de Neyman-Pearson --- Licence}]
\textbf{STA-09.} \textbf{Problème.} On teste une hypothèse nulle simple $ H_0 $ : les données ont pour densité
$ f_0 $, contre une alternative simple $ H_1 $ : densité $ f_1 $.
\begin{enumerate}
\item[(a)] Démontrez le lemme de Neyman-Pearson : parmi tous les tests de risque de première espèce au
plus $ \alpha $, un test qui rejette exactement lorsque $ f_1(x)/f_0(x) >k $ (avec $ k $
choisi pour que le niveau soit $ \alpha $, en randomisant à la frontière si nécessaire) maximise
la puissance $ P_1(\text{rejet}) $.
\item[(b)] Appliquez-le : pour $ n $ observations indépendantes $ N(\theta, 1) $, en testant
$ \theta = 0 $ contre $ \theta = \theta_1 $ avec $ \theta_1 >0 $ fixé, montrez que le test
du rapport de vraisemblance est précisément le test z unilatéral \og{} rejeter lorsque
$ \sqrt{n}\, \bar{x} >z_{\alpha} $ \fg{}, et observez que le même test est le plus puissant pour
tout $ \theta_1 >0 $ simultanément (un test uniformément le plus puissant).
\end{enumerate}

\end{problem}

\noindent L'article de 1933 de Neyman et Pearson a transformé le test d'hypothèse d'un
ensemble d'habitudes en un problème d'optimisation : fixez l'erreur que vous redoutez le plus, puis
minimisez l'autre. Le lemme est le théorème fondamental de ce point de vue, et sa démonstration
tient en deux lignes d'idée enveloppées de comptabilité --- le rapport de vraisemblance ordonne les
points de l'échantillon selon leur rendement probant. Tout, de la théorie de la détection radar à
la conception des essais cliniques, repose dessus.

\begin{refs}
\item Contexte : Lemme de Neyman-Pearson --- Wikipédia (\url{https://fr.wikipedia.org/wiki/Lemme_de_Neyman-Pearson})
\item Lecture : E. L. Lehmann \& J. P. Romano, \textit{Testing Statistical Hypotheses}, ch. 3
\item Article : J. Neyman \& E. S. Pearson, \og{} On the problem of the most efficient tests of statistical hypotheses \fg{}, \textit{Philosophical Transactions of the Royal Society A} 231 (1933), 289--337
\end{refs}

\bigskip

\begin{problem}[{Lois a priori conjuguées et mise à jour bayésienne --- Licence}]
\textbf{STA-10.} \textbf{Problème.}
\begin{enumerate}
\item[(a)] Supposons que la probabilité de succès $ p $ d'une pièce ait une loi a priori
$ \mathrm{Beta}(a, b) $ et que l'on observe $ n $ lancers indépendants comportant $ s $
succès. Établissez la loi a posteriori et montrez qu'il s'agit de
$ \mathrm{Beta}(a + s,\, b + n - s) $.
\item[(b)] Montrez que l'espérance a posteriori est une moyenne pondérée de l'espérance a priori
$ a/(a+b) $ et de la proportion observée $ s/n $, et identifiez les poids ; interprétez
$ a + b $ comme une \og{} taille d'échantillon a priori \fg{}.
\item[(c)] Établissez les couples conjugués analogues : loi a priori gamma avec des données de Poisson, et
loi a priori normale avec des données normales de variance connue.
\item[(d)] Démontrez que la mise à jour bayésienne est cohérente : mettre à jour d'un coup sur tout le jeu
de données donne la même loi a posteriori que mettre à jour séquentiellement, une observation à la
fois.
\end{enumerate}

\end{problem}

\noindent Le calcul bêta-Bernoulli est essentiellement celui de l'essai posthume de
Thomas Bayes de 1763, et Laplace utilisa la même machinerie pour estimer la probabilité du lever du
soleil et la masse de Saturne. L'expression \og{} loi a priori conjuguée \fg{} vient de Raiffa et Schlaifer
(1961). La conjugaison est ce qui a rendu l'inférence bayésienne calculable pendant deux siècles,
avant les méthodes de Monte-Carlo par chaînes de Markov ; la partie (d) est l'expression
mathématique de la promesse bayésienne selon laquelle l'ordre dans lequel vous apprenez des faits
ne peut pas compter.

\begin{refs}
\item Contexte : Loi a priori conjuguée --- Wikipédia (en anglais) (\url{https://en.wikipedia.org/wiki/Conjugate_prior})
\item Lecture : A. Gelman et al., \textit{Bayesian Data Analysis}, 3e éd., ch. 2
\item Lecture : L. Wasserman, \textit{All of Statistics}, ch. 11
\end{refs}

\bigskip

\begin{problem}[{Le problème du char d'assaut allemand --- Licence}]
\textbf{STA-11.} \textbf{Problème.} Les numéros de série $ 1, 2, \dots, N $ sont estampillés sur $ N $ chars ;
vous en capturez $ k $, tirés sans remise, tous les sous-ensembles étant équiprobables, et vous
observez des numéros de série de maximum $ M $.
\begin{enumerate}
\item[(a)] Calculez la loi de $ M $ et montrez que l'estimateur du maximum de vraisemblance de $ N $
est $ M $ lui-même --- et qu'il est biaisé par défaut, toujours.
\item[(b)] Démontrez que $ \hat{N} = M\left(1 + \frac{1}{k}\right) - 1 $ est un estimateur sans biais de
$ N $ (\og{} ajoutez l'écart moyen \fg{}).
\item[(c)] Comparez $ \hat{N} $ à l'estimateur des moments $ 2\bar{x} - 1 $ construit sur la moyenne
empirique : calculez les deux variances (ou argumentez laquelle est la plus petite), et expliquez
pourquoi le maximum est ici un meilleur résumé de l'échantillon que la moyenne.
\end{enumerate}

\end{problem}

\noindent Pendant la Seconde Guerre mondiale, les statisticiens alliés estimèrent la
production allemande de chars à partir des numéros de série des chars capturés et détruits. Leur
estimation pour la période juin 1940 -- septembre 1942 tournait autour de 246 chars par mois, tandis
que le renseignement classique annonçait environ 1400 ; les registres de production capturés
montrèrent ensuite un chiffre d'environ 245. L'épisode, documenté par Ruggles et Brodie en 1947,
est la démonstration la plus nette qui soit qu'un bon estimateur bat un grand réseau d'espionnage.
Le problème est aussi un superbe pendant en population finie de l'estimation de la borne d'une loi
uniforme.

\begin{refs}
\item Contexte : Problème du char d'assaut allemand --- Wikipédia (\url{https://fr.wikipedia.org/wiki/Probl%C3%A8me_du_char_d'assaut_allemand})
\item Article : R. Ruggles \& H. Brodie, \og{} An empirical approach to economic intelligence in World War II \fg{}, \textit{Journal of the American Statistical Association} 42 (1947), 72--91
\item Cours : MIT OCW 18.650 Statistics for Applications (\url{https://ocw.mit.edu/courses/18-650-statistics-for-applications-fall-2016/})
\end{refs}

\bigskip

\begin{problem}[{Égaliser les moments --- Licence, short}]
\textbf{STA-22.} \textbf{Problème.} Soient $ x_1, \dots, x_n $ des tirages indépendants dans la
famille gamma de densité
$ f(x) = \frac{\beta^{\alpha}}{\Gamma(\alpha)} x^{\alpha-1} e^{-\beta x} $ sur $ x >0 $,
d'espérance $ \alpha/\beta $ et de variance $ \alpha/\beta^2 $. Établissez les estimateurs des
moments de $ \alpha $ et $ \beta $ en égalant les deux premiers moments empiriques à leurs
valeurs théoriques, et démontrez que les deux estimateurs convergent en probabilité vers les vrais
paramètres, en nommant les deux théorèmes que vous utilisez.
\end{problem}

\noindent Karl Pearson introduisit la méthode dans son mémoire de 1894 sur l'évolution, où
il ajustait un mélange de deux lois normales à des mesures de crabes de Naples et dut résoudre une
équation de degré neuf à la main ; ce fut la première procédure d'estimation à usage général de la
statistique. Fisher montra plus tard que le maximum de vraisemblance (STA-05) est asymptotiquement
plus efficace, et le répéta souvent sans ménagement --- mais les estimateurs des moments survivent
parce qu'ils n'exigent aucune vraisemblance, fournissent d'excellents points de départ pour une
maximisation numérique et, sous la forme généralisée due à Hansen (1982), sont l'outil standard de
l'économétrie.

\begin{refs}
\item Contexte : Méthode des moments (statistiques) --- Wikipédia (\url{https://fr.wikipedia.org/wiki/M%C3%A9thode_des_moments_(statistiques)})
\item Article : K. Pearson, \og{} Contributions to the mathematical theory of evolution \fg{}, \textit{Philosophical Transactions of the Royal Society A} 185 (1894), 71--110
\item Lecture : G. Casella \& R. L. Berger, \textit{Statistical Inference}, 2e éd., ch. 7
\end{refs}

\bigskip

\begin{problem}[{Le coût des questions multiples --- Licence, short}]
\textbf{STA-23.} \textbf{Problème.} Vous testez $ m $ hypothèses nulles et rejetez $ H_i $
dès que sa valeur p vérifie $ p_i \le \alpha/m $. Démontrez que le taux d'erreur par famille --- la
probabilité d'au moins un rejet à tort --- est au plus $ \alpha $ quelle que soit la dépendance
entre les tests, et démontrez que lorsque les $ m $ hypothèses nulles sont toutes vraies et que
les valeurs p sont indépendantes et uniformes, ce taux vaut exactement
$ 1 - (1 - \alpha/m)^m $ ; quantifiez de combien cela reste en dessous de $ \alpha $ pour
$ m $ grand.
\end{problem}

\noindent Toute la démonstration se réduit à l'inégalité de Boole, ce qui explique que la
correction soit indépassable sans hypothèses supplémentaires et, en même temps, bien trop prudente
lorsque les tests sont positivement liés. Elle porte le nom de Carlo Emilio Bonferroni pour ses
inégalités de 1936, mais son usage en comparaisons multiples est dû à Olive Jean Dunn (1961), qui
en mérite davantage le crédit. Les études d'association pangénomique testent environ un million de
marqueurs et déclarent donc la significativité à $ 5 \times 10^{-8} $ ; le fait que ce seuil
écarte de vraies découvertes est ce qui a motivé le taux de fausses découvertes de STA-17(b).

\begin{refs}
\item Contexte : Correction de Bonferroni --- Wikipédia (\url{https://fr.wikipedia.org/wiki/Correction_de_Bonferroni})
\item Contexte : Taux d'erreur par famille --- Wikipédia (en anglais) (\url{https://en.wikipedia.org/wiki/Family-wise_error_rate})
\item Article : O. J. Dunn, \og{} Multiple comparisons among means \fg{}, \textit{Journal of the American Statistical Association} 56 (1961), 52--64
\end{refs}

\bigskip

\begin{problem}[{La pente est-elle réelle ? --- Licence}]
\textbf{STA-24.} \textbf{Problème.} Prenez le modèle de régression simple normal : les $ x_i $ sont fixés et non
tous égaux, et
\[ y_i = a + b x_i + \varepsilon_i , \qquad \varepsilon_1, \dots, \varepsilon_n
 \ \text{indépendants} \ N(0, \sigma^2), \qquad n \ge 3 . \]
Notons $ S_{xx} = \sum_i (x_i - \bar{x})^2 $, soient $ \hat{a} $ et $ \hat{b} $ l'ordonnée à
l'origine et la pente des moindres carrés de STA-04, et soit
$ \mathrm{RSS} = \sum_i (y_i - \hat{a} - \hat{b}x_i)^2 $.
\begin{enumerate}
\item[(a)] Démontrez que $ \hat{b} $ est une fonction linéaire des $ y_i $, donc de loi normale, avec
$ \mathbb{E}[\hat{b}] = b $ et $ \operatorname{Var}(\hat{b}) = \sigma^2 / S_{xx} $.
\item[(b)] Démontrez que $ \mathrm{RSS}/\sigma^2 $ suit la loi du khi-deux à $ n-2 $ degrés de liberté
et est indépendante de $ \hat{b} $. (Décomposez $ \mathbb{R}^n $ en l'espace des colonnes de
dimension deux de la matrice de plan et son orthogonal, et utilisez l'invariance par rotation de la
loi normale multivariée --- cette indépendance est un miracle purement gaussien, et c'est le cœur du
problème.)
\item[(c)] Déduisez-en qu'avec $ \hat{\sigma}^2 = \mathrm{RSS}/(n-2) $,
\[ T = \frac{\hat{b} - b}{\hat{\sigma}/\sqrt{S_{xx}}} \]
suit la loi de Student à $ n-2 $ degrés de liberté, et écrivez le test de niveau $ \alpha $ de
$ H_0: b = 0 $ qui en résulte ainsi que l'intervalle de confiance à $ 95\% $ pour
$ b $.
\item[(d)] Démontrez que la valeur observée de $ T $ sous $ H_0: b = 0 $ vaut
$ r\sqrt{n-2}\,/\sqrt{1-r^2} $, où $ r $ est la corrélation empirique des $ x $ et des
$ y $. Expliquez ensuite, en un paragraphe soigneux et avec STA-03 en tête, ce que le rejet de
$ H_0 $ vous autorise exactement à dire.
\end{enumerate}

\end{problem}

\noindent William Sealy Gosset, brasseur chez Guinness travaillant sur des échantillons de
quatre ou cinq, trouva en 1908 la loi qui permet d'utiliser un écart type estimé à la place d'un
écart type connu ; Guinness interdisant à ses employés de publier, il signa \og{} Student \fg{}. Fisher
fournit la géométrie rigoureuse en dimension $ n $ --- la décomposition orthogonale de la partie
(b) --- et l'étendit à la régression, où elle est devenue l'énoncé inférentiel le plus utilisé de
toute la science : presque toute affirmation publiée de la forme \og{} l'effet était significatif
(p < 0,05) \fg{} est le calcul que l'on vous demande ici de justifier. La partie (d) rappelle que la
mathématique certifie une pente, pas une cause.

\begin{refs}
\item Contexte : Loi de Student --- Wikipédia (\url{https://fr.wikipedia.org/wiki/Loi_de_Student})
\item Contexte : Régression linéaire simple --- Wikipédia (en anglais) (\url{https://en.wikipedia.org/wiki/Simple_linear_regression})
\item Article : Student (W. S. Gosset), \og{} The probable error of a mean \fg{}, \textit{Biometrika} 6 (1908), 1--25
\item Lecture : G. Casella \& R. L. Berger, \textit{Statistical Inference}, 2e éd., ch. 11
\end{refs}

\bigskip

\begin{problem}[{Une seule table pour toutes les lois --- Licence, short}]
\textbf{STA-33.} \textbf{Problème.} Soient $ X_1, \dots, X_n $ indépendantes de fonction de
répartition continue $ F $, soit $ \hat{F}_n(x) = \frac{1}{n}\#\{ i : X_i \le x \} $ la
fonction de répartition empirique, et soit
$ D_n = \sup_{x \in \mathbb{R}} \big| \hat{F}_n(x) - F(x) \big| $. Démontrez que les variables
$ U_i = F(X_i) $ sont indépendantes et uniformes sur $ (0,1) $, déduisez-en que la loi de
$ D_n $ est la même pour toute $ F $ continue, et montrez sur un exemple que la conclusion tombe
en défaut lorsque $ F $ a un atome.
\end{problem}

\noindent C'est pour cela que le test de Kolmogorov-Smirnov a pu exister dès 1933 : une
seule table de valeurs critiques, calculée une fois pour toutes, sert pour toute hypothèse nulle
continue, parce que la transformation intégrale de probabilité efface tout ce qui concerne $ F $
hormis sa continuité. Kolmogorov trouva cette année-là la loi limite de $ \sqrt{n}\,D_n $ et
Smirnov la version à deux échantillons en 1939 ; Doob expliqua en 1949 que la limite est en réalité
une fonctionnelle du pont brownien, ce que Donsker rendit rigoureux en 1952 --- le processus
empirique, l'un des piliers de la statistique non paramétrique moderne. Le pendant à échantillon
fini est l'inégalité de Dvoretzky-Kiefer-Wolfowitz (1956), dont la constante optimale fut fournie
par Massart en 1990, et qui transforme $ D_n $ en une bande de confiance honnête pour une
fonction de répartition inconnue.

\begin{refs}
\item Contexte : Test de Kolmogorov-Smirnov --- Wikipédia (\url{https://fr.wikipedia.org/wiki/Test_de_Kolmogorov-Smirnov})
\item Contexte : Transformation intégrale de probabilité --- Wikipédia (en anglais) (\url{https://en.wikipedia.org/wiki/Probability_integral_transform})
\item Lecture : A. W. van der Vaart, \textit{Asymptotic Statistics}, ch. 19
\end{refs}

\bigskip

\begin{problem}[{Les colliders : la corrélation que l'ajustement fabrique --- Licence}]
\textbf{STA-34.} \textbf{Problème.} Un \textit{collider} est une variable influencée par deux autres. Conditionner
par un collider est l'image en miroir du conditionnement par un facteur de confusion (STA-03) : au
lieu de supprimer une association fallacieuse, il en fabrique une.
\begin{enumerate}
\item[(a)] La version de Berkson. Soient $ A $ et $ B $ deux événements indépendants avec
$ 0 < P(A), P(B) < 1 $ --- deux maladies sans lien --- et soit $ H = A \cup B $ l'événement
\og{} la personne est hospitalisée \fg{}, qui se produit si l'une ou l'autre maladie est présente.
Démontrez que
\[ P(A \mid B \cap H) \;<\; P(A \mid H) , \]
de sorte qu'à l'intérieur de l'hôpital les deux maladies paraissent négativement associées alors
que dans la population elles sont indépendantes.
\item[(b)] La version gaussienne. Soient $ X $ et $ Y $ des variables normales standard indépendantes
et soit $ Z = X + Y + \varepsilon $ avec $ \varepsilon \sim N(0, \sigma^2) $ indépendante des
deux. Démontrez que la corrélation partielle de $ X $ et $ Y $ sachant $ Z $ vaut
\[ \rho_{XY \cdot Z} = -\frac{1}{1 + \sigma^2} , \]
de sorte que \og{} contrôler par $ Z $ \fg{} produit une forte association négative entre deux variables
qui sont exactement indépendantes.
\item[(c)] La version régression. Dans le modèle de (b), démontrez que le coefficient théorique des
moindres carrés porté par $ Y $ dans la régression de $ X $ sur $ (Y, Z) $ est non nul, alors
que le coefficient porté par $ Y $ dans la régression de $ X $ sur $ Y $ seul est nul.
Comparez avec le modèle de confusion $ X = Z + U $, $ Y = Z + V $ de STA-03(a), où ajouter
$ Z $ à la régression est exactement ce qu'il faut faire, et énoncez la règle générale sur les
variables que l'on peut légitimement ajuster.
\item[(d)] Quelle part de tout cela les données seules peuvent-elles révéler ? Considérez les trois
diagrammes causaux sur $ (X, Y, Z) $ : la chaîne $ X \to Z \to Y $, la fourche
$ X \leftarrow Z \rightarrow Y $ et le collider $ X \to Z \leftarrow Y $. Démontrez que les deux
premiers impliquent $ X \perp Y \mid Z $ tandis que le troisième implique $ X \perp Y $ et
généralement pas $ X \perp Y \mid Z $. Concluez que le collider est détectable à partir de la loi
conjointe des observables mais qu'un facteur de confusion ne peut pas être distingué d'un
médiateur --- de sorte qu'ajuster par un médiateur détruit un effet réel exactement comme ajuster par
un collider en invente un faux, et qu'aucune quantité de données ne vous dira lequel des deux vous
êtes en train de faire.
\end{enumerate}

\end{problem}

\noindent Joseph Berkson, médecin et statisticien à la Mayo Clinic, publia la partie (a)
en 1946 pour expliquer pourquoi les associations entre maladies relevées dans les dossiers
hospitaliers ne sont pas fiables, et David Sackett confirma plus tard l'effet empiriquement sous le
nom de \og{} biais d'admission \fg{}. Judea Pearl a donné au phénomène son nom moderne et son critère
graphique, et la partie (d) est la raison pour laquelle ses diagrammes causaux ne sont pas
décoratifs : les indépendances conditionnelles que les données peuvent révéler ne déterminent pas
la structure causale, de sorte que le choix des variables à ajuster est une hypothèse, formulée
avant l'analyse. L'enjeu n'est pas académique --- on a montré en 2020 que les premières études sur la
gravité du COVID-19 fondées sur des patients hospitalisés ou testés étaient déformées par exactement
ce mécanisme, et le fameux \og{} paradoxe de l'obésité \fg{} en cardiologie est largement tenu pour un biais
de collider déguisé.

\begin{refs}
\item Contexte : Paradoxe de Berkson --- Wikipédia (en anglais) (\url{https://en.wikipedia.org/wiki/Berkson%27s_paradox})
\item Contexte : Collisionneur (statistiques) --- Wikipédia (\url{https://fr.wikipedia.org/wiki/Collisionneur_(statistiques)})
\item Article : J. Berkson, \og{} Limitations of the application of fourfold table analysis to hospital data \fg{}, \textit{Biometrics Bulletin} 2 (1946), 47--53
\item Article : G. Griffith et al., \og{} Collider bias undermines our understanding of COVID-19 disease risk and severity \fg{}, \textit{Nature Communications} 11 (2020), 5749
\end{refs}

\bigskip


\section*{Master}

\begin{problem}[{Pourquoi le bootstrap fonctionne --- et quand il échoue --- Master}]
\textbf{STA-12.} \textbf{Problème.} Soient $ x_1, \dots, x_n $ des tirages indépendants d'une loi $ F $, et
soit $ \hat{F}_n $ la loi empirique qui met la masse $ 1/n $ sur chaque observation. Le
bootstrap estime la loi d'échantillonnage d'une statistique sous $ F $ par sa loi
d'échantillonnage sous $ \hat{F}_n $, que l'on peut simuler par rééchantillonnage.
\begin{enumerate}
\item[(a)] Démontrez que pour chaque $ t $ fixé, $ \hat{F}_n(t) \to F(t) $ presque sûrement, et
énoncez le théorème de Glivenko-Cantelli qui rend cette convergence uniforme en $ t $. Expliquez,
en un paragraphe précis, comment le \og{} principe de substitution \fg{} (plug-in) transforme cela en
justification heuristique du bootstrap.
\item[(b)] Cassez-le maintenant. Soit $ F $ la loi uniforme sur $ (0, \theta) $ et prenons pour
statistique le maximum de l'échantillon $ M_n $. Montrez que
$ P(M_n^* = M_n) = 1 - (1 - 1/n)^n \to 1 - 1/e \approx 0{,}632 $, où $ M_n^* $ est le maximum
d'un rééchantillon bootstrap --- de sorte que la loi bootstrap du maximum (renormalisé) a un atome de
masse environ 0,632 en zéro, alors que la vraie loi limite de $ n(\theta - M_n)/\theta $ est une
loi exponentielle continue. Concluez que le bootstrap est inconsistant pour l'extrême.
\item[(c)] Expliquez, au niveau des idées, quelle caractéristique des statistiques \og{} régulières \fg{} comme
les moyennes et les quantiles fait défaut au maximum.
\end{enumerate}

\end{problem}

\noindent Bradley Efron a introduit le bootstrap dans un article des \textit{Annals of
Statistics} de 1979 au titre désarmant de modestie, et il est devenu l'un des outils les plus
utilisés de l'âge informatique : un substitut universel aux calculs analytiques de théorie de
l'échantillonnage. Mais c'est un théorème, pas de la magie --- les démonstrations de consistance
(Bickel-Freedman, Singh, 1981) ont des hypothèses, et le maximum de l'échantillon est le
contre-exemple standard qui montre qu'elles sont nécessaires. Comprendre les deux moitiés de ce
problème, c'est comprendre le bootstrap.

\begin{refs}
\item Contexte : Bootstrap (statistiques) --- Wikipédia (\url{https://fr.wikipedia.org/wiki/Bootstrap_(statistiques)})
\item Lecture : B. Efron \& T. Hastie, \textit{Computer Age Statistical Inference}, ch. 10--11 (PDF gratuit sur le site du livre (\url{https://hastie.su.domains/CASI/}))
\item Lecture : L. Wasserman, \textit{All of Statistics}, ch. 8
\item Article : B. Efron, \og{} Bootstrap methods: Another look at the jackknife \fg{}, \textit{Annals of Statistics} 7 (1979), 1--26
\end{refs}

\bigskip

\begin{problem}[{Exhaustivité et Rao-Blackwell --- Master}]
\textbf{STA-13.} \textbf{Problème.} Une statistique $ T $ est exhaustive pour un paramètre $ \theta $ si la loi
conditionnelle des données sachant $ T $ ne dépend pas de $ \theta $.
\begin{enumerate}
\item[(a)] Démontrez le théorème de factorisation dans le cas discret : $ T $ est exhaustive si et
seulement si la vraisemblance se factorise en $ g(T(x), \theta)\, h(x) $. Utilisez-le pour
montrer que $ \sum_i x_i $ est exhaustive pour des échantillons de Bernoulli et de Poisson.
\item[(b)] Démontrez le théorème de Rao-Blackwell : si $ W $ est un estimateur sans biais d'une quantité
$ \tau(\theta) $ et si $ T $ est exhaustive, alors $ W' = \mathbb{E}[W \mid T] $ est une
véritable statistique (elle ne dépend pas de $ \theta $), est sans biais, et a une variance au
plus égale à celle de $ W $ --- strictement plus petite sauf si $ W $ était déjà une fonction de
$ T $.
\item[(c)] Faites tourner la machine : pour des données de Poisson($ \lambda $), partez de l'estimateur
sans biais grossier $ \mathbf{1}\{x_1 = 0\} $ de $ e^{-\lambda} = P(X = 0) $, conditionnez par
$ T = \sum_i x_i $, et établissez l'estimateur amélioré
$ \left(\frac{n-1}{n}\right)^{T} $. Vérifiez directement qu'il est sans biais.
\end{enumerate}

\end{problem}

\noindent L'exhaustivité est l'idée la plus profonde de Fisher (1920--22) : une statistique
qui extrait des données jusqu'à la dernière goutte d'information sur $\theta$. Rao (1945) et Blackwell
(1947) en ont fait indépendamment un dispositif d'amélioration --- conditionnez n'importe quel
estimateur honnête par une statistique exhaustive et vous ne pouvez que gagner. La partie (c) est
le tour de passe-passe classique : un estimateur qui ne regarde que la première observation est
transformé, par pur conditionnement, en un excellent estimateur. Combinée à la complétude
(Lehmann-Scheffé), cette machine produit les meilleurs estimateurs sans biais.

\begin{refs}
\item Contexte : Théorème de Rao-Blackwell --- Wikipédia (\url{https://fr.wikipedia.org/wiki/Th%C3%A9or%C3%A8me_de_Rao-Blackwell})
\item Contexte : Statistique exhaustive --- Wikipédia (\url{https://fr.wikipedia.org/wiki/Statistique_exhaustive})
\item Lecture : G. Casella \& R. L. Berger, \textit{Statistical Inference}, 2e éd., ch. 6--7
\end{refs}

\bigskip

\begin{problem}[{Information de Fisher et borne de Cramér-Rao --- Master}]
\textbf{STA-14.} \textbf{Problème.} Pour un modèle statistique régulier de densité $ f(x; \theta) $, le score est
la dérivée de la log-vraisemblance en $ \theta $, et l'information de Fisher $ I(\theta) $ est
la variance du score.
\begin{enumerate}
\item[(a)] Sous les conditions de régularité usuelles (dérivation sous le signe intégral), démontrez que
le score est d'espérance nulle et que $ I(\theta) $ est aussi égale à l'opposé de l'espérance de
la dérivée seconde de la log-vraisemblance.
\item[(b)] Démontrez l'inégalité de Cramér-Rao : tout estimateur sans biais $ W $ de $ \theta $ fondé
sur $ n $ observations indépendantes vérifie
$ \operatorname{Var}(W) \ge \frac{1}{n\, I(\theta)} $. (Le cœur de la démonstration est
l'inégalité de Cauchy-Schwarz appliquée à la covariance de $ W $ avec le score.)
(c) Calculez $ I(\theta) $ pour les modèles Bernoulli($ p $), Poisson($ \lambda $) et
$ N(\mu, \sigma^2) $ à $ \sigma $ connu, et exhibez dans chaque cas un estimateur sans biais
atteignant la borne.
(d) Expliquez (démonstration non exigée) en quel sens l'EMV est asymptotiquement normal de variance
$ \frac{1}{n\, I(\theta)} $ --- la justification ultime de STA-05.
\end{enumerate}

\end{problem}

\noindent L'information de Fisher, introduite dans les articles de Fisher de 1922--25, est
la courbure de la vraisemblance : des vraisemblances pointues épinglent le paramètre. La borne
inférieure fut démontrée par C. R. Rao (1945) et figure dans le traité de Harald Cramér de 1946 ;
elle marque la frontière exacte de ce que l'estimation sans biais peut atteindre, et l'information
de Fisher traverse aujourd'hui la statistique, la théorie de l'information et jusqu'à la géométrie
des familles de lois (géométrie de l'information). Gardez ce problème en tête lorsque vous
rencontrerez STA-15, qui montre ce qui arrive quand on abandonne l'absence de biais.

\begin{refs}
\item Contexte : Borne de Cramér-Rao --- Wikipédia (\url{https://fr.wikipedia.org/wiki/Borne_de_Cram%C3%A9r-Rao})
\item Contexte : Information de Fisher --- Wikipédia (\url{https://fr.wikipedia.org/wiki/Information_de_Fisher})
\item Lecture : G. Casella \& R. L. Berger, \textit{Statistical Inference}, 2e éd., ch. 7 et 10
\end{refs}

\bigskip

\begin{problem}[{Le paradoxe de James-Stein --- Master}]
\textbf{STA-15.} \textbf{Problème.} On observe un unique $ X = (X_1, \dots, X_p) $ à
composantes indépendantes $ X_i \sim N(\theta_i, 1) $, et l'on estime le vecteur entier
$ \theta $ sous la perte quadratique totale. L'estimateur évident est $ X $ lui-même, de risque
constant $ p $. L'estimateur de James-Stein rétrécit vers l'origine :
\[ \delta(X) = \left( 1 - \frac{p-2}{\lVert X \rVert^2} \right) X ,
\qquad \lVert X \rVert^2 = \sum_{i=1}^{p} X_i^2 . \]
(a) Démontrez l'identité de Stein : pour $ g $ suffisamment régulière avec
$ \mathbb{E}\left| \partial g/\partial x_i \right| $ finie,
$ \mathbb{E}[(X_i - \theta_i)\, g(X)] = \mathbb{E}[\partial g/\partial x_i] $.
(b) Utilisez-la pour calculer le risque de $ \delta $ et démontrez que pour $ p \ge 3 $,
\[ \mathbb{E}\,\lVert \delta(X) - \theta \rVert^2
 = p - (p-2)^2\, \mathbb{E}\!\left[ \frac{1}{\lVert X \rVert^2} \right] <p
\quad \text{pour tout } \theta \]
--- l'estimateur naturel est inadmissible en dimension trois ou plus.
(c) Expliquez pourquoi l'argument échoue pour $ p = 2 $, et méditez le scandale : estimer la
moyenne au bâton d'un joueur en utilisant les moyennes de joueurs sans rapport diminue l'erreur
totale. Où exactement l'\og{} emprunt de force \fg{} entre-t-il dans la mathématique ?
\end{problem}

\noindent Charles Stein annonça en 1956 que le vecteur des moyennes empiriques est
inadmissible en dimension $\ge$ 3, et avec Willard James exhiba en 1961 l'estimateur dominant explicite
ci-dessus. Cela reste le théorème le plus contre-intuitif de la statistique : ni l'absence de biais
ni aucun principe d'optimalité n'y survit sans réserve. L'article d'Efron et Morris de 1977 dans
\textit{Scientific American} l'a rendu célèbre avec des données de baseball, et ses descendants
intellectuels --- bayésien empirique, rétrécissement, régression ridge, modèles hiérarchiques --- sont
l'ossature de l'estimation moderne à grande échelle.

\begin{refs}
\item Contexte : Estimateur de James-Stein --- Wikipédia (en anglais) (\url{https://en.wikipedia.org/wiki/James%E2%80%93Stein_estimator})
\item Lecture : B. Efron \& T. Hastie, \textit{Computer Age Statistical Inference}, ch. 7 (PDF gratuit sur le site du livre (\url{https://hastie.su.domains/CASI/}))
\item Article : B. Efron \& C. Morris, \og{} Stein's paradox in statistics \fg{}, \textit{Scientific American} 236 (1977), 119--127
\item Article : W. James \& C. Stein, \og{} Estimation with quadratic loss \fg{}, \textit{Proc. Fourth Berkeley Symposium} (1961), 361--379
\end{refs}

\bigskip

\begin{problem}[{La méthode delta --- Master, short}]
\textbf{STA-25.} \textbf{Problème.} Supposons que $ \sqrt{n}\,(T_n - \theta) $ converge en loi
vers $ N(0, \sigma^2) $ et que $ g $ soit dérivable en $ \theta $ avec
$ g'(\theta) \ne 0 $. Démontrez que
\[ \sqrt{n}\,\big( g(T_n) - g(\theta) \big) \ \xrightarrow{\ d\ } \
 N\!\big(0,\; g'(\theta)^2 \sigma^2\big) , \]
en précisant explicitement où le théorème de Slutsky est utilisé, et appliquez-le pour trouver la
loi limite de la cote empirique $ \hat{p}/(1-\hat{p}) $ pour des données de Bernoulli($ p $)
indépendantes avec $ 0 <p <1 $.
\end{problem}

\noindent C'est le théorème de Taylor faisant de la statistique : une fonction régulière
d'un estimateur asymptotiquement normal est asymptotiquement normale, avec la variance multipliée
par le carré de la dérivée. Les physiciens du XIXe siècle connaissaient cela sous le nom
de propagation des erreurs ; Cramér en donna l'énoncé moderne en 1946. Presque toute erreur type
que vous lirez sur une sortie d'ordinateur pour un rapport, un log-odds, une corrélation ou une
probabilité de survie a été produite par cet argument --- et lorsque $ g'(\theta) = 0 $ il échoue
et la limite devient une loi du khi-deux, ce qui vaut la peine d'être découvert par soi-même.

\begin{refs}
\item Contexte : Méthode delta --- Wikipédia (\url{https://fr.wikipedia.org/wiki/M%C3%A9thode_delta})
\item Lecture : A. W. van der Vaart, \textit{Asymptotic Statistics}, ch. 3
\item Lecture : G. Casella \& R. L. Berger, \textit{Statistical Inference}, 2e éd., ch. 5
\end{refs}

\bigskip

\begin{problem}[{Biais et variance, en une seule identité --- Master, short}]
\textbf{STA-26.} \textbf{Problème.} Soit $ Y = f(x) + \varepsilon $ en un point fixé $ x $,
avec $ \mathbb{E}[\varepsilon] = 0 $ et $ \operatorname{Var}(\varepsilon) = \sigma^2 $, et soit
$ \hat{f} $ une estimation de $ f $ construite à partir d'un échantillon d'apprentissage
indépendant de $ \varepsilon $. Démontrez la décomposition
\[ \mathbb{E}\big[(Y - \hat{f}(x))^2\big] = \sigma^2
 + \big( \mathbb{E}[\hat{f}(x)] - f(x) \big)^2 + \operatorname{Var}\big(\hat{f}(x)\big) , \]
et servez-vous-en pour expliquer précisément pourquoi un estimateur biaisé peut avoir une erreur de
prédiction plus petite que tous les estimateurs sans biais.
\end{problem}

\noindent Trois lignes de calcul, et elles réorganisent toute une discipline : le premier
terme est irréductible, le deuxième est ce que vous achetez en simplifiant votre modèle, le
troisième est ce que vous payez pour la souplesse. L'identité est la justification honnête des
estimateurs à rétrécissement de STA-15 et de toute pénalité de régularisation en apprentissage
automatique ; Geman, Bienenstock et Doursat ont baptisé le compromis correspondant \og{} dilemme
biais-variance \fg{} en 1992. Elle a aussi une variante moderne qui mérite lecture : les modèles très
surparamétrés présentent une seconde descente de l'erreur de test au-delà du point d'interpolation,
de sorte que l'image classique en U est un cas particulier plutôt qu'une loi.

\begin{refs}
\item Contexte : Dilemme biais-variance --- Wikipédia (\url{https://fr.wikipedia.org/wiki/Dilemme_biais-variance})
\item Article : S. Geman, E. Bienenstock \& R. Doursat, \og{} Neural networks and the bias/variance dilemma \fg{}, \textit{Neural Computation} 4 (1992), 1--58
\item Lecture : T. Hastie, R. Tibshirani \& J. Friedman, \textit{The Elements of Statistical Learning}, ch. 7 (PDF gratuit sur le site du livre (\url{https://hastie.su.domains/ElemStatLearn/}))
\end{refs}

\bigskip

\begin{problem}[{Ce que mesure vraiment une courbe ROC --- Master}]
\textbf{STA-27.} \textbf{Problème.} Un détecteur attribue à chaque objet un score réel $ s(X) $. Soit $ X_1 $
un objet tiré au hasard dans la classe positive et $ X_0 $ un objet de la classe négative, et
pour un seuil $ t $ posons
\[ \mathrm{TPR}(t) = P\big(s(X_1) \ge t\big), \qquad
 \mathrm{FPR}(t) = P\big(s(X_0) \ge t\big) . \]
La courbe ROC est l'ensemble des points $ (\mathrm{FPR}(t), \mathrm{TPR}(t)) $ lorsque $ t $
parcourt $ \mathbb{R} $, et l'AUC est l'aire sous cette courbe.
\begin{enumerate}
\item[(a)] Démontrez que la courbe ROC est inchangée si l'on remplace le score par $ h(s(X)) $ pour
toute $ h $ strictement croissante --- la courbe ROC ne voit donc que l'ordre des objets, jamais la
calibration des scores.
\item[(b)] Démontrez l'identité probabiliste
\[ \mathrm{AUC} = P\big( s(X_1) >s(X_0) \big) + \tfrac12 P\big( s(X_1) = s(X_0) \big) , \]
et identifiez sa version empirique comme la statistique U de Mann-Whitney divisée par le produit
des tailles d'échantillon.
\item[(c)] Supposons que les deux classes aient les densités $ f_1 $ et $ f_0 $. Démontrez que la
courbe ROC du score du rapport de vraisemblance $ s^*(x) = f_1(x)/f_0(x) $ est au-dessus ou
confondue avec la courbe ROC de tout autre score --- c'est le lemme de Neyman-Pearson (STA-09) lu
comme un énoncé sur des courbes --- et démontrez que cette courbe optimale est concave, sa pente au
point de fonctionnement valant le seuil de rapport de vraisemblance correspondant.
\item[(d)] Voici maintenant le piège pratique : démontrez que l'AUC ne change pas lorsque la proportion de
positifs dans la population change, alors que la précision
$ P(\text{positif} \mid \text{signalé}) $ change. Construisez deux problèmes de dépistage ayant
des courbes ROC identiques dont l'un est utile et l'autre déclenche cent fausses alertes par cas
réel, et énoncez la morale quant à la publication d'un chiffre unique.
\end{enumerate}

\end{problem}

\noindent La caractéristique de fonctionnement du récepteur est née du radar : pendant la
Seconde Guerre mondiale, la question était de savoir comment régler le seuil de décision d'un
opérateur quand un écho peut être un avion ou du bruit, et la théorie fut rédigée par Peterson,
Birdsall et Fox en 1954. Green et Swets la portèrent en psychologie expérimentale, et l'article de
Hanley et McNeil de 1982 --- qui a fait connaître l'équivalence de (b) --- l'a rendue standard en
radiologie puis dans tous les domaines qui évaluent des classifieurs. La mathématique de (c) est un
théorème centenaire en costume moderne ; l'avertissement de (d) explique que l'AUC soit à la fois le
nombre le plus rapporté et le plus critiqué de l'apprentissage automatique appliqué.

\begin{refs}
\item Contexte : Courbe ROC --- Wikipédia (\url{https://fr.wikipedia.org/wiki/Courbe_ROC})
\item Contexte : Test de Wilcoxon-Mann-Whitney --- Wikipédia (\url{https://fr.wikipedia.org/wiki/Test_de_Wilcoxon-Mann-Whitney})
\item Article : J. A. Hanley \& B. J. McNeil, \og{} The meaning and use of the area under a receiver operating characteristic (ROC) curve \fg{}, \textit{Radiology} 143 (1982), 29--36
\item Lecture : D. M. Green \& J. A. Swets, \textit{Signal Detection Theory and Psychophysics} (1966)
\end{refs}

\bigskip

\begin{problem}[{Un estimateur minimax pour une pièce --- Master, short}]
\textbf{STA-35.} \textbf{Problème.} Soit $ X \sim \mathrm{Binomial}(n, p) $ et estimons
$ p $ sous la perte quadratique, de sorte que le risque d'un estimateur $ \delta $ vaut
$ R(p, \delta) = E_p\big[ (\delta(X) - p)^2 \big] $. Démontrez qu'un estimateur de Bayes dont la
fonction de risque est constante en $ p $ est minimax, puis démontrez que l'estimateur de Bayes
pour la loi a priori $ \mathrm{Beta}(\sqrt{n}/2, \sqrt{n}/2) $,
\[ \delta^{*}(X) = \frac{X + \sqrt{n}/2}{n + \sqrt{n}} , \]
a un risque constant et est donc minimax ; comparez enfin $ R(p, \delta^{*}) $ au risque de la
proportion empirique $ X/n $ et déterminez pour quels $ p $ l'estimateur ordinaire est
meilleur.
\end{problem}

\noindent Abraham Wald a bâti la statistique comme un jeu à deux joueurs contre la nature
dans les années 1940 (\textit{Statistical Decision Functions}, 1950), et la règle minimax est le coup
du pessimiste : choisir la procédure dont le pire cas est le moins mauvais. Le procédé employé ici ---
trouver une loi a priori dont la règle de Bayes a un risque plat, et le pire cas n'a alors nulle
part où se cacher --- est la voie standard vers la minimaxité, et la loi a priori qui fait cela
s'appelle la moins favorable. Hodges et Lehmann ont traité ce cas et des cas voisins en 1950. La
comparaison finale est la coda honnête : l'estimateur minimax rétrécit vers $ 1/2 $ et ne bat la
proportion empirique qu'au voisinage du milieu de l'intervalle, ce qui rappelle que la minimaxité
est une garantie sur le pire cas et non une recommandation pour tous les cas.

\begin{refs}
\item Contexte : Estimateur minimax --- Wikipédia (en anglais) (\url{https://en.wikipedia.org/wiki/Minimax_estimator})
\item Article : J. L. Hodges \& E. L. Lehmann, \og{} Some problems in minimax point estimation \fg{}, \textit{Annals of Mathematical Statistics} 21 (1950), 182--197
\item Lecture : E. L. Lehmann \& G. Casella, \textit{Theory of Point Estimation}, 2e éd., ch. 5
\end{refs}

\bigskip

\begin{problem}[{Pourquoi l'algorithme EM grimpe --- Master}]
\textbf{STA-36.} \textbf{Problème.} Supposons que les données observées $ y $ soient une fonction de données
complètes $ (y, z) $ où $ z $ n'est pas observé, et écrivons
$ \ell(\theta) = \log p(y \mid \theta) $. L'algorithme EM itère les deux étapes
\[ Q(\theta \mid \theta^{(t)}) = E\big[ \log p(y, z \mid \theta) \,\big|\, y, \theta^{(t)} \big],
 \qquad
 \theta^{(t+1)} = \arg\max_{\theta} \, Q(\theta \mid \theta^{(t)}) . \]
\begin{enumerate}
\item[(a)] Démontrez la propriété de croissance. Montrez que
$ \ell(\theta) = Q(\theta \mid \theta') - H(\theta \mid \theta') $ où
$ H(\theta \mid \theta') = E\big[ \log p(z \mid y, \theta) \mid y, \theta' \big] $, démontrez à
l'aide de l'inégalité de Jensen que $ H(\theta' \mid \theta') - H(\theta \mid \theta') $ est une
divergence de Kullback-Leibler, donc positive, et concluez que
$ \ell(\theta^{(t+1)}) \ge \ell(\theta^{(t)}) $ pour tout $ t $ --- la vraisemblance des données
observées ne décroît jamais, quel que soit le modèle.
\item[(b)] Démontrez l'identité d'énergie libre : pour toute loi $ q $ sur la variable latente,
\[ \ell(\theta) = \underbrace{E_q\big[\log p(y,z \mid \theta)\big] + \mathcal{H}(q)}_{F(q,\theta)}
 \; + \; \mathrm{KL}\big( q \,\big\|\, p(z \mid y, \theta) \big) , \]
où $ \mathcal{H}(q) $ est l'entropie de $ q $, et déduisez-en que EM est exactement une montée
par coordonnées sur $ F $ : l'étape E maximise en $ q $, l'étape M en $ \theta $.
\item[(c)] Traitez un exemple : pour un mélange gaussien univarié à deux composantes d'espérances,
variances et poids de mélange inconnus, établissez les responsabilités de l'étape E et montrez que
l'étape M est un maximum de vraisemblance pondéré ordinaire. Démontrez ensuite que cette
vraisemblance n'est pas bornée --- envoyez l'espérance d'une composante sur un point de données et sa
variance vers zéro --- de sorte que \og{} EM converge vers l'estimateur du maximum de vraisemblance \fg{} ne
peut pas être vrai tel quel.
\item[(d)] Dites précisément ce qui est vrai. Donnez un exemple à plusieurs maxima locaux où EM converge
vers des points différents selon le point de départ, énoncez le théorème de Wu de 1983 sur la
convergence vers un point stationnaire de $ \ell $ sous conditions de régularité, et expliquez ce
que la propriété de croissance de (a) garantit et ne garantit pas au sujet de la suite
$ \theta^{(t)} $ elle-même.
\end{enumerate}

\end{problem}

\noindent Dempster, Laird et Rubin ont nommé l'algorithme et démontré la propriété de
croissance en 1977, unifiant une douzaine de procédures inventées séparément --- celle de Hartley
pour les données groupées (1958), l'algorithme de Baum-Welch pour les modèles de Markov cachés
(1970), l'estimation de fréquences géniques, l'analyse factorielle --- en une seule recette dont le
seul ingrédient est la capacité de calculer une espérance conditionnelle. C'est l'un des articles
les plus cités de la statistique. Son argument de convergence n'était toutefois pas étanche :
C. F. Jeff Wu signala la faille en 1983 et fournit l'énoncé correct, ce qui est l'objet de la partie
(d) --- une suite croissante de valeurs de la vraisemblance ne dit rien à elle seule sur la
convergence des paramètres, et encore moins sur leur convergence vers un maximum global. La partie
(b) est la reformulation de Neal et Hinton de 1998, et c'est la porte d'entrée de l'inférence
variationnelle : restreignez $ q $ à une famille traitable et la même montée par coordonnées
devient le cheval de bataille du calcul bayésien moderne.

\begin{refs}
\item Contexte : Algorithme espérance-maximisation --- Wikipédia (\url{https://fr.wikipedia.org/wiki/Algorithme_esp%C3%A9rance-maximisation})
\item Article : A. P. Dempster, N. M. Laird \& D. B. Rubin, \og{} Maximum likelihood from incomplete data via the EM algorithm \fg{}, \textit{Journal of the Royal Statistical Society B} 39 (1977)
\item Article : C. F. J. Wu, \og{} On the convergence properties of the EM algorithm \fg{}, \textit{Annals of Statistics} 11 (1983)
\item Lecture : G. McLachlan \& T. Krishnan, \textit{The EM Algorithm and Extensions}, 2e éd.
\end{refs}

\bigskip

\begin{problem}[{Survie, censure et une hypothèse que l'on ne peut pas tester --- Master}]
\textbf{STA-37.} \textbf{Problème.} Chaque sujet $ i $ possède un temps de survie $ T_i $ et un temps de censure
$ C_i $ ; on n'observe que $ Y_i = \min(T_i, C_i) $ et l'indicatrice
$ \delta_i = \mathbf{1}\{ T_i \le C_i \} $. Notons $ S(t) = P(T > t) $ la fonction de survie
à estimer.
\begin{enumerate}
\item[(a)] Supposons $ T_i $ et $ C_i $ indépendants, la loi de $ C_i $ ne dépendant pas des
paramètres d'intérêt. Démontrez que la vraisemblance se factorise en
$ \prod_i f(y_i)^{\delta_i} S(y_i)^{1-\delta_i} $ multiplié par un terme ne faisant pas
intervenir $ S $, de sorte que le mécanisme de censure peut être ignoré. Énoncez exactement ce
que la censure indépendante affirme au sujet du monde, et expliquez pourquoi c'est une hypothèse
sur des mécanismes plutôt qu'une propriété des données observées.
\item[(b)] Établissez l'estimateur de Kaplan-Meier
\[ \hat{S}(t) = \prod_{j \,:\, t_j \le t} \left( 1 - \frac{d_j}{n_j} \right) , \]
où $ t_1 < t_2 < \cdots $ sont les instants de défaillance observés distincts, $ d_j $ le
nombre de défaillances en $ t_j $ et $ n_j $ le nombre de sujets encore à risque juste avant,
en maximisant la vraisemblance de (a) sur toutes les lois portées par les instants de défaillance
observés. Établissez ensuite la formule de Greenwood pour la variance de $ \hat{S}(t) $ en
appliquant la méthode delta (STA-25) à $ \log \hat{S}(t) $.
\item[(c)] Le modèle de Cox attribue au sujet $ i $ de vecteur de covariables $ x_i $ le risque
instantané $ \lambda(t \mid x_i) = \lambda_0(t)\, e^{\beta^{\top} x_i} $ où $ \lambda_0 $ est
une fonction inconnue arbitraire. Établissez la vraisemblance partielle
\[ L(\beta) = \prod_{i \,:\, \delta_i = 1}
 \frac{ e^{\beta^{\top} x_i} }{ \sum_{j \in R(t_i)} e^{\beta^{\top} x_j} } , \qquad
 R(t) = \{ j : Y_j \ge t \} , \]
démontrez que $ \lambda_0 $ en a entièrement disparu, et démontrez que son score est d'espérance
nulle en la vraie valeur de $ \beta $. Énoncez ensuite quelle quantité est estimée lorsque
l'hypothèse de risques proportionnels est fausse.
\item[(d)] Démontrez ou reconstruisez le théorème de Tsiatis (1975) : pour toute loi conjointe de
$ (T, C) $ --- avec une dépendance arbitraire --- il existe un modèle où $ T $ et $ C $ sont
\textit{indépendants} qui produit exactement la même loi des observables $ (Y, \delta) $. Concluez
que la censure indépendante n'est pas testable à partir des données, et discutez de ce à quoi
devrait ressembler une analyse de sensibilité pour cette hypothèse.
\end{enumerate}

\end{problem}

\noindent Les tables de mortalité sont plus anciennes que la statistique --- John Graunt a
dressé celle de Londres en 1662 et Halley celle de Breslau en 1693 --- mais la censure, le fait
qu'une étude s'achève alors que des sujets sont encore en vie, fut traitée par des recettes
actuarielles ad hoc jusqu'à ce qu'Edward Kaplan et Paul Meier soumettent indépendamment au
\textit{JASA} des articles presque identiques ; le rédacteur en chef les obligea à fusionner, et après
quatre ans de discussions l'article de 1958 parut. Il est parmi les plus cités de l'histoire des
sciences, tout comme la vraisemblance partielle de David Cox de 1972, si nouvelle que Cox dut écrire
un second article en 1975 pour expliquer pourquoi on peut traiter une vraisemblance partielle comme
une vraisemblance. La partie (d) est l'ombre permanente du domaine : l'hypothèse qui fait tout
fonctionner --- que les gens quittent une étude pour des raisons sans rapport avec leur pronostic --- ne
peut pas être vérifiée sur les données, seulement argumentée.

\begin{refs}
\item Contexte : Estimateur de Kaplan-Meier --- Wikipédia (\url{https://fr.wikipedia.org/wiki/Estimateur_de_Kaplan-Meier})
\item Contexte : Régression de Cox (modèle à risques proportionnels) --- Wikipédia (\url{https://fr.wikipedia.org/wiki/R%C3%A9gression_de_Cox})
\item Article : E. L. Kaplan \& P. Meier, \og{} Nonparametric estimation from incomplete observations \fg{}, \textit{Journal of the American Statistical Association} 53 (1958), 457--481
\item Article : D. R. Cox, \og{} Regression models and life-tables \fg{}, \textit{Journal of the Royal Statistical Society B} 34 (1972), 187--202 ; A. Tsiatis, \og{} A nonidentifiability aspect of the problem of competing risks \fg{}, \textit{PNAS} 72 (1975), 20--22
\end{refs}

\bigskip

\begin{problem}[{AIC et BIC répondent à des questions différentes --- Master}]
\textbf{STA-38.} \textbf{Problème.} Un modèle à $ k $ paramètres libres est ajusté à $ n $ observations
indépendantes, donnant la log-vraisemblance maximisée $ \hat{\ell} $. Les deux critères de
sélection classiques sont
\[ \mathrm{AIC} = -2\hat{\ell} + 2k , \qquad \mathrm{BIC} = -2\hat{\ell} + k \log n . \]
Ils ne diffèrent que par une pénalité, et ils sont dérivés de cibles complètement différentes.
\begin{enumerate}
\item[(a)] Établissez l'AIC. Prenez pour cible la divergence de Kullback-Leibler moyenne de la vérité au
modèle ajusté, autrement dit la log-vraisemblance moyenne du modèle ajusté sur une copie
indépendante des données. Montrez que $ \hat{\ell} $ surestime cette cible, et démontrez sous les
conditions de régularité usuelles --- un développement quadratique de la log-vraisemblance plus la
normalité asymptotique de l'estimateur du maximum de vraisemblance (STA-14) --- que le biais vaut
asymptotiquement $ k $.
\item[(b)] Établissez le BIC. Appliquez l'approximation de Laplace à la vraisemblance marginale
$ m(y) = \int p(y \mid \theta)\, \pi(\theta)\, d\theta $ et démontrez que
\[ \log m(y) = \hat{\ell} - \frac{k}{2}\log n + O(1) , \]
de sorte que comparer des valeurs de BIC approche la comparaison des probabilités a posteriori des
modèles ; montrez que la loi a priori $ \pi $ ne contribue qu'au terme $ O(1) $.
\item[(c)] Voyez la conséquence dans le cas le plus simple. Pour $ n $ observations de
$ N(\mu, 1) $, comparez les deux critères sur le couple emboîté $ \{\mu = 0\} $ contre
$ \{\mu \in \mathbb{R}\} $ : démontrez que lorsque $ \mu = 0 $ la règle du BIC choisit le bon
modèle avec une probabilité tendant vers 1, tandis que la règle de l'AIC choisit le plus grand
modèle avec une probabilité qui ne s'annule pas, et identifiez la valeur limite de cette
probabilité. Démontrez ensuite, ou énoncez fidèlement, le résultat de Stone de 1977 selon lequel
l'AIC est asymptotiquement équivalent à la validation croisée \og{} un contre tous \fg{}.
\item[(d)] Énoncez précisément le théorème de Yang (2005) : aucune procédure de sélection de modèle ne
peut être simultanément consistante pour identifier un vrai modèle de dimension finie et optimale
au sens des vitesses minimax pour estimer la fonction de régression --- les forces du BIC et celles
de l'AIC ne peuvent pas être combinées. Puis, pour un problème appliqué concret de votre choix,
dites lequel des deux objectifs est réellement le vôtre.
\end{enumerate}

\end{problem}

\noindent Hirotugu Akaike proposa son critère lors d'un colloque en 1971 et le publia en
1974, ayant compris que la vraisemblance maximisée est une estimation biaisée de l'ajustement
prédictif et que le biais vaut, chose remarquable, tout juste le nombre de paramètres. Gideon
Schwarz établit l'autre pénalité en 1978 par un argument bayésien où la taille de l'échantillon
entre par le rétrécissement de la loi a posteriori. On s'est ensuite disputé à leur sujet pendant
quarante ans comme si l'un des deux devait avoir tort, mais la partie (d) tranche le débat de la
seule façon dont les mathématiques le peuvent : ils optimisent des critères différents, et le
théorème d'impossibilité de Yuhong Yang démontre qu'aucun critère n'optimise les deux. C'est un bon
modèle de la manière dont s'achève une controverse méthodologique --- non par un vainqueur, mais par
un théorème décrivant l'arbitrage qu'aucun des deux camps ne pouvait éviter.

\begin{refs}
\item Contexte : Critère d'information d'Akaike --- Wikipédia (\url{https://fr.wikipedia.org/wiki/Crit%C3%A8re_d'information_d'Akaike})
\item Contexte : Critère d'information bayésien --- Wikipédia (\url{https://fr.wikipedia.org/wiki/Crit%C3%A8re_d'information_bay%C3%A9sien})
\item Article : H. Akaike, \og{} A new look at the statistical model identification \fg{}, \textit{IEEE Transactions on Automatic Control} 19 (1974), 716--723 ; G. Schwarz, \og{} Estimating the dimension of a model \fg{}, \textit{Annals of Statistics} 6 (1978)
\item Article : Y. Yang, \og{} Can the strengths of AIC and BIC be shared? \fg{}, \textit{Biometrika} 92 (2005), 937--950
\end{refs}

\bigskip


\section*{Recherche}

\begin{problem}[{L'inférence après sélection --- Recherche}]
\textbf{STA-16.} \textbf{Problème.} La pratique moderne choisit la question après avoir vu les données --- le plus
grand effet, le meilleur modèle --- puis rapporte l'inférence comme si la question avait été fixée
d'avance. Rendez les dégâts précis, puis passez les réparations en revue.
\begin{enumerate}
\item[(a)] Soient $ X_1, \dots, X_m $ indépendantes $ N(\mu_i, 1) $ et supposons que tous les
$ \mu_i = 0 $. Vous sélectionnez l'indice du plus grand $ X_i $ et publiez la borne de confiance
inférieure naïve à 95 \% $ X_{(m)} - 1{,}645 $ pour l'espérance sélectionnée. Démontrez que la
couverture vaut $ \Phi(1{,}645)^m $, soit environ 0,90 pour $ m = 2 $ et qui s'effondre vers 0
quand $ m $ croît. Montrez aussi que
$ \mathbb{E}[\max(X_1, X_2)] = 1/\sqrt{\pi} >0 $, de sorte que l'estimation sélectionnée est
biaisée par excès alors même que toutes les vraies espérances sont nulles (la malédiction du
vainqueur).
\item[(b)] Problème de compréhension : lisez comment on restaure une inférence post-sélection valide ---
inférence simultanée sur tous les modèles (Berk et al. 2013) et \og{} inférence sélective \fg{}
conditionnelle exacte pour des procédures comme le lasso (école de Taylor et Tibshirani) --- et
rédigez une comparaison de ce que chacune garantit, et à quel prix en largeur d'intervalle.
\textit{État des lieux : domaine de recherche actif depuis 2013 environ ; aucun cadre unique n'est
encore la pratique standard, et étendre l'inférence sélective exacte au-delà de procédures
particulières reste un territoire ouvert.}
\end{enumerate}

\end{problem}

\noindent Le problème est aussi ancien que le tripatouillage des données, mais il est
devenu aigu quand la sélection de modèle est devenue automatique : le lasso, la régression pas à
pas et les balayages pangénomiques choisissent tous les hypothèses de façon adaptative. L'article
de 2013 de Berk, Brown, Buja, Zhang et Zhao a donné la première réponse fréquentiste pleinement
honnête (\og{} PoSI \fg{}), et l'inférence sélective conditionnelle a suivi rapidement. C'est le domaine où
la mathématique de la crise de la reproductibilité (STA-17) se répare, une procédure à la fois.

\begin{refs}
\item Contexte : Data dredging (triturage de données) --- Wikipédia (\url{https://fr.wikipedia.org/wiki/Data_dredging})
\item Article : R. Berk, L. Brown, A. Buja, K. Zhang \& L. Zhao, \og{} Valid post-selection inference \fg{}, \textit{Annals of Statistics} 41 (2013), 802--837
\item Article : J. Taylor \& R. J. Tibshirani, \og{} Statistical learning and selective inference \fg{}, \textit{PNAS} 112 (2015), 7629--7634
\end{refs}

\bigskip

\begin{problem}[{La crise de la reproductibilité comme statistique mathématique --- Recherche}]
\textbf{STA-17.} \textbf{Problème.} Vers 2011--2015, il est apparu qu'une large part des résultats publiés en
psychologie, en médecine et ailleurs ne se reproduisent pas. Traitez cela comme un domaine où l'on
démontre des théorèmes, pas comme un scandale.
\begin{enumerate}
\item[(a)] Un domaine teste de nombreuses hypothèses, dont une fraction $ \pi $ correspond à de vrais
effets ; les tests ont un niveau $ \alpha $ et une puissance $ 1 - \beta $. Établissez la valeur
prédictive positive
\[ \mathrm{PPV} = \frac{\pi (1-\beta)}{\pi (1-\beta) + (1-\pi)\alpha} , \]
et évaluez-la pour $ \pi = 0{,}1 $, $ \alpha = 0{,}05 $, puissance $ 0{,}8 $. Concluez que des
résultats \og{} significatifs \fg{} peuvent être faux plus d'une fois sur trois même avec une pratique
irréprochable --- c'est l'argument d'Ioannidis de 2005 rendu rigoureux.
\item[(b)] Démontrez que la procédure de Benjamini-Hochberg contrôle le taux de fausses découvertes au
niveau $ q $ lorsque les valeurs p des vraies hypothèses nulles sont indépendantes et
uniformes.
\item[(c)] Démontrez que l'arrêt optionnel détruit la significativité naïve : pour des observations
indépendantes $ N(\theta, 1) $ avec $ \theta = 0 $, la loi du logarithme itéré implique que la
statistique z courante $ \sqrt{n}\, \bar{x}_n $ dépasse tout seuil fixé une infinité de fois,
presque sûrement --- de sorte que \og{} tester après chaque observation et s'arrêter à la significativité \fg{}
rejette une hypothèse nulle vraie avec probabilité 1.
\textit{État des lieux : (a)--(c) sont des mathématiques établies ; concevoir une inférence robuste au
biais de publication, au p-hacking et à l'analyse adaptative est un front de recherche ouvert et
actif (voir STA-16 et STA-18).}
\end{enumerate}

\end{problem}

\noindent L'essai d'Ioannidis de 2005 \og{} Why most published research findings are false \fg{}
et la réplication par l'Open Science Collaboration en 2015 de 100 études de psychologie (39 \% des
résultats significatifs répliqués) ont défini la crise. Les réponses mathématiques --- taux de fausses
découvertes (Benjamini-Hochberg 1995), inférence post-sélection, préenregistrement, méthodes
valides à tout instant --- comptent parmi les statistiques les plus conséquentes que l'on fasse
aujourd'hui : les théorèmes présentés ici sont la raison pour laquelle la crise a pu être
diagnostiquée.

\begin{refs}
\item Contexte : Crise de la reproductibilité --- Wikipédia (\url{https://fr.wikipedia.org/wiki/Crise_de_la_reproductibilit%C3%A9})
\item Article : J. P. A. Ioannidis, \og{} Why most published research findings are false \fg{}, \textit{PLoS Medicine} 2 (2005), e124
\item Article : Y. Benjamini \& Y. Hochberg, \og{} Controlling the false discovery rate \fg{}, \textit{Journal of the Royal Statistical Society B} 57 (1995), 289--300
\item Article : Open Science Collaboration, \og{} Estimating the reproducibility of psychological science \fg{}, \textit{Science} 349 (2015), aac4716
\end{refs}

\bigskip

\begin{problem}[{Valeurs e et tests sûrs --- Recherche}]
\textbf{STA-18.} \textbf{Problème.} Une \textit{variable e} pour une hypothèse nulle $ H_0 $ est une variable
aléatoire positive $ E $ telle que $ \mathbb{E}_0[E] \le 1 $ sous toute loi de $ H_0 $ ; sa
valeur réalisée est une valeur e, les grandes valeurs comptant comme preuve contre
$ H_0 $.
\begin{enumerate}
\item[(a)] Démontrez que rejeter lorsque $ E \ge 1/\alpha $ a un risque de première espèce au plus
$ \alpha $ (inégalité de Markov), et que $ \min(1, 1/E) $ est une valeur p conservatrice.
\item[(b)] Démontrez que pour une hypothèse nulle simple de densité $ f_0 $, tout rapport de
vraisemblance $ f_1(X)/f_0(X) $ est une variable e.
\item[(c)] Voici l'intérêt de tout cela : supposons que des expériences arrivent en séquence et que chaque
nouvelle variable e $ E_k $ puisse être choisie à partir de toutes les données passées, avec
$ \mathbb{E}_0[E_k \mid \text{passé}] \le 1 $. Démontrez que le produit courant
$ M_n = E_1 E_2 \cdots E_n $ est une surmartingale positive sous $ H_0 $, énoncez l'inégalité de
Ville $ P_0\!\left( \sup_n M_n \ge 1/\alpha \right) \le \alpha $, et concluez que la preuve
empirique peut être accumulée et surveillée en continu, avec arrêt optionnel et poursuite
optionnelle, sans rien coûter à la garantie de première espèce --- exactement ce que les valeurs p ne
peuvent pas faire (STA-17c).
\item[(d)] Problème de compréhension : lisez la théorie d'optimalité au taux de croissance (GRO) de
Grünwald, de Heide et Koolen et expliquez ce qui remplace la \og{} puissance \fg{} lorsque les tests se
mesurent à la croissance de la valeur e.
\textit{État des lieux : programme de recherche fondationnel, très actif depuis 2019 ; \og{} Safe testing \fg{}
a paru avec discussion dans JRSS-B en 2024, et les questions ouvertes abondent (valeurs e optimales
pour des hypothèses nulles composites, valeurs e pour des données complexes dépendantes, le bon
analogue des intervalles de confiance).}
\end{enumerate}

\end{problem}

\noindent Les valeurs e renouent le test d'hypothèse avec les martingales et le jeu --- la
preuve contre l'hypothèse nulle est la fortune d'un joueur qui parie contre elle --- une ligne de
pensée qui remonte à Ville (1939) et que la probabilité théorie-des-jeux de Shafer et Vovk a
ravivée. L'essor moderne a commencé avec \og{} Safe testing \fg{} de Grünwald, de Heide et Koolen
(prépublication 2019, JRSS-B 2024) et les travaux parallèles de Vovk, Wang, Ramdas et Shafer ; le
domaine se range désormais sous la bannière \og{} statistique théorie-des-jeux et inférence sûre valide
à tout instant \fg{}. C'est la tentative actuelle la plus sérieuse de refonder le test de sorte que la
pratique scientifique ordinaire --- jeter un œil aux données, les mettre en commun, poursuivre ---
devienne mathématiquement légale.

\begin{refs}
\item Contexte : Valeurs e --- Wikipédia (en anglais) (\url{https://en.wikipedia.org/wiki/E-values})
\item Article : P. Grünwald, R. de Heide \& W. M. Koolen, \og{} Safe testing \fg{}, \textit{Journal of the Royal Statistical Society B} 86 (2024), 1091--1128, arXiv:1906.07801 (\url{https://arxiv.org/abs/1906.07801})
\item Article : A. Ramdas, P. Grünwald, V. Vovk \& G. Shafer, \og{} Game-theoretic statistics and safe anytime-valid inference \fg{}, \textit{Statistical Science} 38 (2023), 576--601, arXiv:2210.01948 (\url{https://arxiv.org/abs/2210.01948})
\item Lecture : G. Shafer \& V. Vovk, \textit{Game-Theoretic Foundations for Probability and Finance} (2019)
\end{refs}

\bigskip

\begin{problem}[{Ce que la prédiction conforme peut et ne peut pas promettre --- Recherche, short}]
\textbf{STA-28.} \textbf{Problème.} La prédiction conforme par découpage ajuste un prédicteur
quelconque sur une partie des données, calcule des scores de conformité $ R_1, \dots, R_n $ (par
exemple des résidus absolus) sur un ensemble de calibration mis de côté, et pour une nouvelle
entrée renvoie l'ensemble des valeurs dont le score tomberait en dessous du
$ \lceil (n+1)(1-\alpha) \rceil $-ième plus petit des $ R_i $. Démontrez la garantie de
couverture marginale à échantillon fini
$ P\big( Y_{n+1} \in C(X_{n+1}) \big) \ge 1 - \alpha $ en supposant seulement que les $ n+1 $
couples sont échangeables --- aucune correction du modèle, aucune hypothèse de loi --- puis énoncez
précisément le théorème d'impossibilité de Barber, Candès, Ramdas et Tibshirani (2021), à savoir
qu'aucune procédure sans hypothèse de loi ne peut atteindre une couverture non triviale
conditionnellement à $ X_{n+1} = x $.
\end{problem}

\noindent La prédiction conforme a été développée par Vovk, Gammerman et Shafer à partir
de la fin des années 1990 et est devenue, autour de 2019, la façon standard d'attacher une barre
d'erreur honnête à la sortie d'un modèle boîte noire arbitraire, y compris des réseaux de neurones
que personne ne sait analyser. La garantie de la première moitié est un argument de permutation en
trois lignes ; le théorème de la seconde moitié explique pourquoi les praticiens continuent de
réclamer quelque chose que le cadre ne peut démontrablement pas leur donner gratuitement. État en
2026 : la couverture marginale sous échangeabilité est une mathématique réglée, tandis que les
relâchements utiles --- couverture conditionnelle approchée, changement de covariables, données
dépendantes et séries temporelles --- forment une frontière de recherche active.

\begin{refs}
\item Contexte : Prédiction conforme --- Wikipédia (en anglais) (\url{https://en.wikipedia.org/wiki/Conformal_prediction})
\item Lecture : V. Vovk, A. Gammerman \& G. Shafer, \textit{Algorithmic Learning in a Random World} (2005)
\item Article : R. F. Barber, E. J. Candès, A. Ramdas \& R. J. Tibshirani, \og{} The limits of distribution-free conditional predictive inference \fg{}, \textit{Information and Inference} 10 (2021), 455--482, arXiv:1903.04684 (\url{https://arxiv.org/abs/1903.04684})
\end{refs}

\bigskip

\begin{problem}[{Un écart entre ce qui peut être estimé et ce qui peut être calculé --- Recherche, short}]
\textbf{STA-29.} \textbf{Problème.} Dans le modèle de covariance à pic, on observe $ n $
tirages indépendants de $ N(0,\, I_p + \theta vv^{\top}) $, où le vecteur unitaire $ v $ a au
plus $ k $ coordonnées non nulles, et la tâche est de décider si $ \theta = 0 $ ou
$ \theta >0 $. Rédigez un compte rendu précis de l'état des lieux : la recherche exhaustive
sur tous les supports réussit dès que $ n \gtrsim k \log p / \theta^2 $, toute méthode connue en
temps polynomial exige $ n \gtrsim k^2 / \theta^2 $, et Berthet et Rigollet (2013) ont démontré
que franchir la seconde frontière résoudrait aussi le problème de la clique plantée --- de sorte que
l'écart apparent n'est réel que conditionnellement à une hypothèse de difficulté en moyenne qui
n'est elle-même pas démontrée.
\end{problem}

\noindent Les écarts calcul-statistique sont l'un des thèmes marquants de la statistique
du XXIe siècle : la théorie classique demande ce qu'une taille d'échantillon rend
possible, et ne dit rien de ce qu'un algorithme peut atteindre en temps polynomial. La détection
d'une composante principale parcimonieuse en est le cas d'école, et les indices en faveur de
l'écart viennent aujourd'hui de plusieurs directions --- réductions depuis la clique plantée, bornes
inférieures somme-de-carrés, et l'heuristique des polynômes de bas degré --- dont aucune n'est une
démonstration inconditionnelle. État en 2026 : l'écart est conjectural ; démontrer
inconditionnellement qu'un problème d'estimation naturel est difficile en calcul reste hors de
portée.

\begin{refs}
\item Contexte : ACP parcimonieuse --- Wikipédia (en anglais) (\url{https://en.wikipedia.org/wiki/Sparse_PCA})
\item Article : Q. Berthet \& P. Rigollet, \og{} Complexity theoretic lower bounds for sparse principal component detection \fg{}, \textit{Proceedings of COLT 2013}, PMLR 30, 1046--1066, PMLR (\url{https://proceedings.mlr.press/v30/Berthet13.html})
\end{refs}

\bigskip

\begin{problem}[{Quand un nombre mérite-t-il le nom d'effet ? --- Recherche}]
\textbf{STA-30.} \textbf{Problème.} Dans le modèle des résultats potentiels, chaque unité $ i $ porte deux
nombres $ Y_i(1) $ et $ Y_i(0) $ --- sa réponse sous traitement et sous contrôle --- dont seul
$ Y_i = Y_i(T_i) $ est observé, où $ T_i \in \{0,1\} $ enregistre le traitement effectivement
reçu. L'estimande est l'effet moyen du traitement
$ \tau = \mathbb{E}[Y(1) - Y(0)] $.
\begin{enumerate}
\item[(a)] Démontrez que $ \tau $ n'est \textit{pas} déterminé par la loi conjointe des observables
$ (T, Y) $ : exhibez deux populations de mêmes lois observables et de $ \tau $ différents.
Démontrez ensuite que sous randomisation --- $ T $ indépendant de $ (Y(0), Y(1)) $ --- la différence
des deux moyennes empiriques est sans biais pour $ \tau $.
\item[(b)] Supposez plutôt l'absence de confusion sachant des covariables $ X $, c'est-à-dire
$ T \perp (Y(0), Y(1)) \mid X $, avec la positivité
$ 0 <P(T=1 \mid X) <1 $ presque sûrement. Démontrez la formule d'ajustement
\[ \mathbb{E}[Y(1)] = \mathbb{E}_X\big[\, \mathbb{E}[Y \mid T=1, X] \,\big] , \]
et démontrez le théorème du score de propension de Rosenbaum et Rubin : l'absence de confusion
sachant $ X $ implique l'absence de confusion sachant le seul nombre
$ e(X) = P(T = 1 \mid X) $.
\item[(c)] Énoncez le critère de la porte dérobée (back-door) de Pearl pour un graphe causal et démontrez
qu'un ensemble le vérifiant justifie la formule d'ajustement de (b). Énoncez ensuite --- sans
démonstration --- le théorème de complétude du do-calcul (Shpitser-Pearl et Huang-Valtorta, 2006) :
dès qu'une loi interventionnelle est identifiable à partir du graphe et de la loi observationnelle,
le do-calcul permet de la dériver.
\item[(d)] Lisez la frontière et rendez-en compte honnêtement : l'analyse de sensibilité (quelle doit être
la force d'un facteur de confusion non mesuré pour renverser une conclusion ?), l'interférence entre
unités (le traitement d'une unité affecte le résultat d'une autre, de sorte que la notation en
résultats potentiels ci-dessus est déjà fausse), et l'estimation de $ \tau $ lorsque les
fonctions de nuisance sont ajustées par apprentissage automatique (apprentissage automatique
double/débiaisé). Pour chacun, dites précisément quelles parties sont des théorèmes et lesquelles
sont ouvertes.
\textit{État des lieux : (a)--(c) sont des mathématiques établies datant de 1974--2006 ; les trois sujets
de (d) sont des domaines de recherche actifs, avec des résultats partiels et aucune théorie générale
établie.}
\end{enumerate}

\end{problem}

\noindent Jerzy Neyman écrivit les résultats potentiels pour des expériences agricoles en
1923 ; Donald Rubin les fit revivre et les généralisa à partir de 1974, et Judea Pearl construisit
le pendant graphique, ce qui lui valut le prix Turing en 2011. Les deux écoles se sont disputées
pendant des décennies et sont aujourd'hui comprises comme des formalismes équivalents pour un même
contenu. Ce qui en fait un problème de recherche plutôt qu'un exercice de manuel, c'est que les
hypothèses qui font tout le travail --- pas de confusion non mesurée, pas d'interférence --- ne sont pas
testables à partir des données, de sorte que la mathématique de l'inférence causale est largement
la mathématique de l'énoncé exact de ce que l'on doit supposer, et de la mesure du tort que l'on
subit si l'on se trompe.

\begin{refs}
\item Contexte : Modèle causal de Neyman-Rubin --- Wikipédia (\url{https://fr.wikipedia.org/wiki/Mod%C3%A8le_causal_de_Neyman-Rubin})
\item Contexte : Inférence causale --- Wikipédia (\url{https://fr.wikipedia.org/wiki/Inf%C3%A9rence_causale})
\item Article : P. R. Rosenbaum \& D. B. Rubin, \og{} The central role of the propensity score in observational studies for causal effects \fg{}, \textit{Biometrika} 70 (1983), 41--55
\item Lecture : J. Pearl, \textit{Causality}, 2e éd. (2009) ; G. Imbens \& D. Rubin, \textit{Causal Inference for Statistics, Social, and Biomedical Sciences} (2015)
\end{refs}

\bigskip

\begin{problem}[{Quand interpoler le bruit est sans danger --- Recherche, short}]
\textbf{STA-39.} \textbf{Problème.} En régression linéaire avec plus de prédicteurs que
d'observations, la solution des moindres carrés de norme $ \ell_2 $ minimale ajuste exactement les
données d'apprentissage --- tous les résidus nuls, bruit compris --- et la théorie classique (STA-26)
prédit le désastre, alors qu'en pratique de tels estimateurs prédisent souvent bien. Rédigez un
compte rendu précis de l'état des lieux : énoncez la caractérisation de Bartlett, Long, Lugosi et
Tsigler (2020) des situations où l'interpolation est bénigne, en termes de rangs effectifs du
spectre de la covariance, énoncez exactement ce que la courbe de double descente de Belkin, Hsu, Ma
et Mandal (2019) démontre et ce qu'elle ne démontre pas, et dites honnêtement quelles parties du
tableau sont des théorèmes sur les modèles linéaires et à noyau et lesquelles restent conjecturales
pour les réseaux profonds.
\end{problem}

\noindent Pendant cinquante ans, la courbe en U de STA-26 a été enseignée comme une loi :
au-delà d'une certaine souplesse, l'erreur de test doit remonter. Vers 2018, les praticiens n'ont
cessé d'exhiber des modèles ayant plus de paramètres que de points de données, qui interpolaient
l'ensemble d'apprentissage et généralisaient quand même, et la théorie a dû être reconstruite plutôt
que défendue. Pour la régression linéaire de norme minimale, la réponse est aujourd'hui
essentiellement complète et quantitative --- le surajustement bénin se produit exactement lorsque le
spectre de la covariance possède beaucoup de directions petites pour absorber le bruit tandis que
quelques grandes portent le signal --- et des résultats analogues existent pour la régression à noyau
sans crête. État en 2026 : réglé pour la régression linéaire et à noyau de norme minimale sous des
conditions spectrales explicites ; pour les réseaux profonds entraînés par descente de gradient, le
phénomène est documenté empiriquement et expliqué seulement dans des régimes restreints.

\begin{refs}
\item Contexte : Double descente --- Wikipédia (en anglais) (\url{https://en.wikipedia.org/wiki/Double_descent})
\item Article : P. L. Bartlett, P. M. Long, G. Lugosi \& A. Tsigler, \og{} Benign overfitting in linear regression \fg{}, \textit{PNAS} 117 (2020), 30063--30070, arXiv:1906.11300 (\url{https://arxiv.org/abs/1906.11300})
\item Article : M. Belkin, D. Hsu, S. Ma \& S. Mandal, \og{} Reconciling modern machine-learning practice and the classical bias--variance trade-off \fg{}, \textit{PNAS} 116 (2019), 15849--15854
\end{refs}

\bigskip

\begin{problem}[{Une inférence qui ne révèle rien sur personne --- Recherche}]
\textbf{STA-40.} \textbf{Problème.} Un algorithme randomisé $ M $ est $ \varepsilon $-\textit{différentiellement
confidentiel} si pour tout couple de jeux de données $ D, D' $ différant par l'enregistrement
d'un seul individu et tout ensemble mesurable $ S $,
\[ P\big( M(D) \in S \big) \;\le\; e^{\varepsilon}\, P\big( M(D') \in S \big) . \]
La question de ce problème est de savoir ce qu'une telle contrainte coûte en précision
statistique.
\begin{enumerate}
\item[(a)] Démontrez la boîte à outils de base. Pour une statistique $ f $ de sensibilité globale
$ \Delta f = \max_{D \sim D'} | f(D) - f(D') | $, démontrez que publier
$ f(D) + \mathrm{Laplace}(\Delta f / \varepsilon) $ est
$ \varepsilon $-différentiellement confidentiel ; démontrez que toute fonction d'une sortie
confidentielle est confidentielle avec le même $ \varepsilon $ (post-traitement) ; et démontrez
qu'exécuter des mécanismes de paramètres $ \varepsilon_1, \dots, \varepsilon_k $ sur les mêmes
données est $ (\varepsilon_1 + \cdots + \varepsilon_k) $-confidentiel.
\item[(b)] Chiffrez la garantie dans le modèle central, où un dépositaire de confiance détient les
données. Pour la moyenne de $ n $ nombres de $ [0,1] $, calculez la sensibilité, et démontrez
que l'écart type du bruit de confidentialité est d'ordre $ 1/(n\varepsilon) $ tandis que l'écart
type d'échantillonnage est d'ordre $ 1/\sqrt{n} $. Concluez précisément pour quelles suites
$ \varepsilon = \varepsilon_n $ le bruit de confidentialité est asymptotiquement
négligeable.
\item[(c)] Chiffrez-la dans le modèle local, où chaque individu randomise avant de publier quoi que ce
soit. Montrez que la réponse randomisée de Warner (1965) --- répondre sincèrement avec probabilité
$ p $, sinon inverser la réponse --- est $ \varepsilon $-confidentielle avec
$ \varepsilon = \log\frac{p}{1-p} $, construisez l'estimateur sans biais de la proportion dans la
population, et démontrez que sa variance est gonflée d'un facteur d'ordre $ 1/\varepsilon^2 $
pour $ \varepsilon $ petit. Énoncez ensuite le théorème minimax de Duchi, Jordan et Wainwright :
sous confidentialité locale, la taille d'échantillon effective est d'ordre $ n\varepsilon^2 $
plutôt que $ n $, et aucune procédure n'y échappe.
\item[(d)] Rendez compte de la frontière en séparant les théorèmes des problèmes ouverts. Considérez
(i) les intervalles de confiance et les tests d'hypothèse à validité exacte en échantillon fini
sous confidentialité, (ii) le choix de $ \varepsilon $, qui est une décision politique sans
aucune théorie statistique pour la guider --- le Bureau du recensement des États-Unis a dû en fixer
un pour le recensement de 2020 et s'est fait attaquer en justice pour les conséquences --- et
(iii) la confidentialité pour des données dont les enregistrements sont dépendants, comme dans les
réseaux ou en cas de participation répétée, où la définition même d'un \og{} jeu de données voisin \fg{}
est contestée. Pour chacun, dites ce qui est démontré et ce qui ne l'est pas.
\textit{État des lieux : les mécanismes de (a) et les vitesses minimax de (b)--(c) sont des mathématiques
établies ; l'inférence confidentielle, la sémantique du budget de confidentialité sur la durée de
vie d'un jeu de données et la confidentialité sous dépendance sont des recherches ouvertes et
actives.}
\end{enumerate}

\end{problem}

\noindent La confidentialité différentielle a été définie par Cynthia Dwork, Frank
McSherry, Kobbi Nissim et Adam Smith en 2006 pour sortir d'un piège : toute notion antérieure
d'anonymisation avait été brisée par des attaques de réidentification, et le domaine avait besoin
d'une garantie ne dépendant pas de ce qu'un attaquant savait déjà. La définition offre exactement
cela --- une borne sur le déplacement que la participation d'une seule personne peut imprimer à la loi
de la sortie --- et elle a été déployée par Apple, par Google et, de la façon la plus lourde de
conséquences, par le recensement américain de 2020. Ce déploiement est le lieu où statistique et
politique publique sont entrées en collision : les démographes objectaient que des effectifs bruités
à petite échelle déforment précisément les populations que le recensement existe pour compter,
tandis que le Bureau soutenait que les attaques par reconstruction sur les données de 2010 avaient
déjà rendu illusoires les anciennes protections. La mathématique de l'arbitrage --- combien de
précision une promesse de confidentialité doit coûter --- est l'un des domaines les plus actifs de la
statistique théorique.

\begin{refs}
\item Contexte : Confidentialité différentielle --- Wikipédia (\url{https://fr.wikipedia.org/wiki/Confidentialit%C3%A9_diff%C3%A9rentielle})
\item Article : S. L. Warner, \og{} Randomized response: a survey technique for eliminating evasive answer bias \fg{}, \textit{Journal of the American Statistical Association} 60 (1965), 63--69
\item Article : J. C. Duchi, M. I. Jordan \& M. J. Wainwright, \og{} Minimax optimal procedures for locally private estimation \fg{} (2016--17), arXiv:1604.02390 (\url{https://arxiv.org/abs/1604.02390})
\item Lecture : C. Dwork \& A. Roth, \textit{The Algorithmic Foundations of Differential Privacy} (2014)
\end{refs}

\bigskip


\section*{Pour aller plus loin}


\vfill
\noindent\emph{Hors de la Caverne} --- des probl\`emes, pas de r\'eponses.
\end{document}
